\documentclass[11pt,a4paper]{article}
\usepackage{fontspec}
\usepackage{polyglossia}
\setdefaultlanguage{russian}
\setotherlanguage{english}
\setmainfont{Liberation Serif}
\setsansfont{DejaVu Sans}
\setmonofont{DejaVu Sans Mono}[Scale=0.85]
\usepackage{amsmath,amssymb,amsthm}
\usepackage[margin=2.3cm]{geometry}
\usepackage{xcolor}
\usepackage{graphicx,booktabs,array,multirow}
\usepackage[font=small,labelfont=bf]{caption}
\usepackage{enumitem}
\usepackage[section]{placeins}
\usepackage[colorlinks=true,linkcolor=blue!50!black,citecolor=blue!50!black,urlcolor=blue!50!black]{hyperref}
\setlist{nosep,leftmargin=1.4em}
\linespread{1.08}
\emergencystretch=1.5em

\theoremstyle{plain}
\newtheorem{theorem}{Теорема}
\newtheorem{lemma}{Лемма}
\newtheorem{proposition}{Предложение}
\newtheorem{corollary}{Следствие}
\newtheorem{conjecture}{Гипотеза}
\theoremstyle{definition}
\newtheorem{definition}{Определение}
\newtheorem{remark}{Замечание}
\newtheorem{claim}{Утверждение}
\renewcommand{\proofname}{Доказательство}

\newcommand{\G}{\Gamma}
\newcommand{\dd}{\,\mathrm{d}}
\DeclareMathOperator{\sign}{sign}

\title{Барьеры порядка явных дробных методов Рунге--Кутты\\ в форме Вольтерры}
\author{Д.\,А.~Приказчиков\\[2pt]
\small Московский физико-технический институт (национальный исследовательский университет),\\
\small г.~Долгопрудный, Московская обл., 141701, Россия\\
\small E-mail: \texttt{prikazchikov.da@phystech.edu}}
\date{Черновик от 16 октября 2026 г.}

\begin{document}
\maketitle

\begin{abstract}
В работе рассматриваются явные одношаговые методы типа Рунге--Кутты для задачи Коши с дробной производной Капуто порядка $\alpha\in(0,1)$. Уравнения такого вида возникают, в частности, при локализации интегродифференциальных уравнений теории вязкоупругости с ядрами типа Работнова, где переменные памяти среды удовлетворяют дробно-дифференциальным уравнениям по времени. Для широко цитируемой конструкции дробных методов Рунге--Кутты с переносом базовой точки производной (EFORK) доказано, что на внутренних шагах локальная ошибка равна $h\,y'(t_n)\bigl(R(\alpha)-1\bigr)+O(h^{1+\alpha})$, где $R(\alpha)=\sum_i w_ic_i^{1-\alpha}/\G(2-\alpha)$. Условие $R(\alpha)=1$ необходимо для согласованности; из условий порядка этих методов оно вытекает лишь при $\alpha\in\{1/2,\dots,1/p\}$~--- именно при таких значениях проводились исходные численные эксперименты, а при прочих $\alpha$ условия порядка его не обеспечивают. Для опубликованных методов при прочих $\alpha$ оно нарушено. Оценка сохраняется и при вычислении интеграла памяти кусочно-линейной интерполяцией, как в опубликованной реализации. Для методов в форме Вольтерры (класс Пуазе--Любича) получено явное представление младших членов локальной ошибки через квадратурные и <<цепочечные>> дефекты и доказан дробный аналог первого барьера Бутчера: наилучший показатель локальной ошибки $O(h^q)$ равен $q=1+2\alpha$ при двух стадиях, $q=\max(2+\alpha,\,1+3\alpha)$ при трёх и не превосходит $\max(2+2\alpha,\,1+4\alpha)$ при четырёх стадиях; при $\alpha=1$ эти формулы переходят в классическое ограничение $p\le s$, при $\alpha=1/2$ согласуются с результатом Любича. Доказательства основаны на обобщённом правиле знаков Декарта для мер, сосредоточенных в узлах метода. Достижимость оценки для четырёх стадий установлена численно в двенадцати значениях $\alpha$. Для пяти стадий численные расчёты приводят к гипотезе о трёхветвевом барьере $q=\max\bigl(2+3\alpha,\,\min(1+5\alpha,\,3+2\alpha)\bigr)$, который при $\alpha\to1$ непрерывно переходит в классическое ограничение $p\le4$. Численно показано также, что вычисление интеграла памяти по значениям стадий снижает глобальный порядок до $\approx1+\alpha$, и предложен способ обойти это ограничение. Построены оптимальные методы с тремя и четырьмя стадиями; их показатели подтверждены прямым измерением локальной ошибки и расчётами на линейной, нелинейной и неавтономной системной задачах. Код и данные открыты.

\medskip\noindent\textbf{Ключевые слова:} дробные дифференциальные уравнения, производная Капуто, методы Рунге--Кутты, уравнения Вольтерры, условия порядка, барьеры Бутчера, вязкоупругость, оператор Работнова.
\end{abstract}

\section{Введение}

Учёт наследственных свойств материалов приобретает большую значимость при моделировании конструкций, работающих в условиях длительных нагрузок. Согласно линейной теории Вольтерры~\cite{Volterra1982}, развитой в работах Работнова~\cite{Rabotnov1966,Rabotnov1977}, определяющие соотношения таких сред содержат интегральные операторы со слабосингулярными ядрами, а естественным математическим аппаратом для их описания служит дробное интегродифференцирование~\cite{SKM1987,Uchaikin2008,Diethelm2010}. В работах~\cite{PPKh2024,PP2026zhvm,Prik2026de} для уравнений линейной вязкоупругости с ядром Работнова была построена система первого порядка, в которой переменные памяти среды $\xi$ удовлетворяют уравнениям вида $D^{\gamma}\xi=-\beta\xi-\beta g(t,x)$, $\gamma\in(0,1)$, т.\,е. задаче Коши для дробно-дифференциального уравнения по времени (в обозначениях~\cite{PP2026zhvm,Prik2026de} $\gamma=-\alpha$, где $\alpha\in(-1,0)$~--- параметр ядра Работнова; при нулевых начальных данных производные Римана--Лиувилля и Капуто совпадают); параметры ядра Работнова для конкретных материалов определяются по кривым ползучести~\cite{PP2025ljm}, а нелинейные обобщения с повреждениями рассмотрены в~\cite{Prik2026sibzhim}. В сеточно-характеристической схеме~\cite{PP2026zhvm} для таких уравнений использовалась разностная аппроксимация невысокого порядка по времени, и повышение этого порядка представляется естественным следующим шагом. Наиболее привлекательны здесь явные одношаговые методы типа Рунге--Кутты: они не требуют решения нелинейных уравнений на шаге и легко встраиваются в явные схемы по пространству.

Итак, рассмотрим задачу Коши
\begin{equation}\label{eq:fde}
{}^{C}\!D^{\alpha}_{0}y(t)=f\bigl(t,y(t)\bigr),\qquad y(0)=y_0,\qquad 0<\alpha<1.
\end{equation}
Её решение совпадает с решением интегрального уравнения Вольтерры
\begin{equation}\label{eq:volterra}
y(t)=y_0+\frac{1}{\G(\alpha)}\int_0^t (t-s)^{\alpha-1}f\bigl(s,y(s)\bigr)\dd s ,
\end{equation}
и основная трудность численного решения связана с нелокальностью: на каждом шаге необходимо учитывать всю предысторию процесса. Среди численных методов выделяются многошаговые схемы и свёрточные квадратуры~\cite{Lubich1985,LubichOstermann1993,BanjaiLopez2018}, схемы предиктор--корректор~\cite{DFF2002}, коллокационные методы~\cite{Brunner2004} и одношаговые методы типа Рунге--Кутты.

Для классических методов Рунге--Кутты теория Бутчера~\cite{Butcher2016,HNW1993} описывает условия порядка через корневые деревья и даёт, в частности, \emph{первый барьер}: явный $s$-стадийный метод имеет порядок не выше $s$. Перенос этой теории на дробный случай возможен двумя способами.

\textbf{Перенос базовой точки.} В работе~\cite{Ghaffari2023} предложены дробные методы Рунге--Кутты (FORK/\allowbreak EFORK), в~которых на каждом шаге базовая точка производной Капуто переносится в $t_n$, правая часть модифицируется приближённым интегралом памяти, а условия порядка $p\alpha$ выводятся из обобщённой формулы Тейлора~\cite{OdibatShawagfeh2007} по степеням $(t-t_n)^{\alpha}$. Численные эксперименты в~\cite{Ghaffari2023} проводились при $\alpha\in\{1/2,1/3,1/4\}$.

\textbf{Форма Вольтерры.} Любич~\cite{Lubich1983} построил теорию методов Рунге--Кутты для уравнений~\eqref{eq:volterra} со слабосингулярным ядром, включая деревья Вольтерры и, для $\alpha=1/2$, барьеры для явных методов (предложение~12 в~\cite{Lubich1983}). Гаррапа~\cite{Garrappa2010} описал алгоритм автоматической генерации условий порядка для таких методов; классический случай гладких ядер разобран в~\cite{BHN1982}. Разграничим известное и новое. Теория Любича и алгоритм Гаррапы позволяют для заданных $s$, $\alpha$ и желаемого показателя выписать систему условий порядка, но не отвечают на вопрос, совместна ли она: барьер~--- это утверждение о несовместности системы при \emph{всех} значениях коэффициентов, и для его доказательства нужна структура самих условий. Генератор условий порядка в настоящей работе (п.~\ref{sec:generator}) является независимой реализацией той же идеи и новым алгоритмом не считается; используется он для проверки и для численного исследования. Новыми являются: явное представление младших членов локальной ошибки через квадратурные дефекты и <<цепочки>>, доказательство барьеров при $s\le4$ для всех $\alpha\in(0,1)$ с помощью правила знаков Декарта и численная карта барьеров при $s=5$. Общий вид барьеров при произвольном $\alpha\in(0,1)$, насколько известно автору, ранее получен не был.

Настоящая работа посвящена именно этому вопросу. Сформулируем основные результаты, сразу указывая их статус. Строго доказаны: необходимое условие согласованности методов EFORK на внутренних шагах и его нарушение для опубликованных методов~\cite{Ghaffari2023} при $\alpha\ne1/m$ (теорема~\ref{thm:efork}); структура младших членов локальной ошибки методов в форме Вольтерры (предложение~\ref{prop:structure}); дробный аналог первого барьера Бутчера для двух и трёх стадий и верхняя оценка для четырёх стадий (теоремы~\ref{thm:s2}--\ref{thm:s4}), причём при четырёх стадиях барьер, как оказалось, порождается членом, связанным с изменением якобиана вдоль решения. Численно установлены: достижимость верхней оценки для четырёх стадий в двенадцати значениях $\alpha$, карта барьеров для $s\le5$, существование патологической ветви пятистадийных методов с неограниченно растущими коэффициентами, снижение глобального порядка до $\approx1+\alpha$ при вычислении интеграла памяти по значениям стадий, показатели построенных методов. Сформулированы как гипотезы: точное значение барьера для четырёх стадий при всех $\alpha$ и трёхветвевой барьер для пяти стадий. Код генератора условий порядка, все расчётные скрипты и данные открыты (см. п.~\ref{sec:generator}).

\section{Математическая постановка и теоремы}

\subsection{Класс методов}

Пусть $t_n=nh$. Разобьём интеграл в~\eqref{eq:volterra} при $t>t_n$ на <<прошлое>> $[0,t_n]$ и <<текущий шаг>> $[t_n,t]$:
\begin{equation}\label{eq:split}
y(t)=\Phi_n(t)+\frac{1}{\G(\alpha)}\int_{t_n}^{t}(t-s)^{\alpha-1}f\bigl(s,y(s)\bigr)\dd s,\qquad
\Phi_n(t)=y_0+\frac{1}{\G(\alpha)}\int_0^{t_n}(t-s)^{\alpha-1}f\bigl(s,y(s)\bigr)\dd s .
\end{equation}
Функцию $\Phi_n$ будем называть \emph{лаг-членом}, или интегралом памяти. Она определяется только пройденной предысторией, но вычисляется в новой точке $t$; поскольку ядро $(t-s)^{\alpha-1}$ зависит от $t$, значения памяти в разных точках текущего шага различны. После замены $s=t_n+\theta h$ интеграл по текущему шагу принимает вид $\frac{h^{\alpha}}{\G(\alpha)}\int_0^{c}(c-\theta)^{\alpha-1}f(t_n+\theta h,\,y(t_n+\theta h))\dd\theta$ при $t=t_n+ch$; отсюда множитель $h^{\alpha}$ в формулах метода.

Явный $s$-стадийный метод в форме Вольтерры (типа Пуазе~\cite{Pouzet1960}) приближает этот интеграл стадиями:
\begin{equation}\label{eq:method}
Y_i=\Phi_n(t_n+c_ih)+h^{\alpha}\sum_{j<i}a_{ij}F_j,\qquad F_j=f(t_n+c_jh,Y_j),\qquad
y_{n+1}=\Phi_n(t_{n+1})+h^{\alpha}\sum_{j=1}^{s}b_jF_j .
\end{equation}
Всюду далее предполагается
\begin{equation}\label{eq:rowsum}
c_1=0,\qquad 0\le c_i\le\hat c,\qquad \sum_{j<i}a_{ij}=\frac{c_i^{\alpha}}{\G(\alpha+1)}\quad(i=2,\dots,s).
\end{equation}
Последнее условие означает, что каждая стадия точно интегрирует постоянные: при $f\equiv1$ точное значение интеграла по текущему шагу равно $\frac{1}{\G(\alpha)}\int_0^{c_i}(c_i-\theta)^{\alpha-1}\dd\theta=\frac{c_i^{\alpha}}{\G(\alpha+1)}$. Это дробный аналог классического условия $\sum_j a_{ij}=c_i$ (условие (9) в~\cite{Lubich1983}). Ограничение $c_i\le\hat c$ естественно: стадия с большим $c_i$ требует значения лаг-члена далеко за пределами текущего шага. В расчётах принято $\hat c=1.5$. Доказательства верхних оценок при $s\le4$ этим ограничением не пользуются, а все найденные оптимальные методы с $s\le4$ имеют $c_i\in[0,1.1]$; существенным оно становится лишь при $s=5$ (п.~\ref{sec:s5}). Обозначим $A=(a_{ij})$, $b=(b_j)$, $c=(c_j)$; степени векторов понимаются покомпонентно.

\paragraph{Классический случай $\alpha=1$.} При $\alpha=1$ ядро тождественно равно единице, уравнение~\eqref{eq:volterra} превращается в интегральную форму задачи Коши $y(t)=y_0+\int_0^tf(s,y(s))\dd s$, используемую в доказательстве теоремы Пикара, а лаг-член перестаёт зависеть от $t$:
\[
\Phi_n(t)=y_0+\int_0^{t_n}f\bigl(s,y(s)\bigr)\dd s=y(t_n).
\]
Вся информация о прошлом сжимается в одно число $y_n$, условие~\eqref{eq:rowsum} переходит в $\sum_ja_{ij}=c_i$, и метод~\eqref{eq:method} становится в точности классическим явным методом Рунге--Кутты:
\[
Y_i=y_n+h\sum_{j<i}a_{ij}F_j,\qquad y_{n+1}=y_n+h\sum_{j=1}^{s}b_jF_j .
\]
Именно поэтому полученные ниже барьеры~--- не аналогия с классикой, а то же утверждение, в которое $\alpha$ входит как параметр. Главное отличие дробного случая состоит в том, что при $\alpha<1$ значения $\Phi_n(t_n+c_ih)$ для разных стадий различны: ядро <<помнит>>, в какую точку текущего шага направлен взгляд.

\begin{definition}
Пусть $n\ge1$, решение $y$ гладкое на $[t_n,t_{n+1}]$, $f$ гладкая. \emph{Локальной ошибкой} называется $\delta_n=y_{n+1}-y(t_{n+1})$, где $y_{n+1}$ вычислено по~\eqref{eq:method} с точным лаг-членом $\Phi_n$. Отметим, что в форме Вольтерры значение $y_n$ в шаг не входит: вся предыстория, включая начальное значение, передаётся через $\Phi_n$, так что точный лаг-член означает и точную предысторию. \emph{Показателем} метода называется наибольшее $q$, при котором $\delta_n=O(h^q)$ для всех гладких $f$ (для систем любой размерности). Через $q(s,\alpha)$ обозначим верхнюю грань показателей по всем явным $s$-стадийным методам класса~\eqref{eq:rowsum}. При $\alpha=1$ показатель равен $p+1$, где $p$~--- классический порядок.
\end{definition}

Подчеркнём, что $q$~--- характеристика \emph{локальной} ошибки одного шага, а не порядок сходимости. Для классического одношагового метода локальная ошибка $O(h^{p+1})$ отвечает глобальному порядку $p$, поскольку ошибки шагов складываются. В рассматриваемых схемах связь иная и зависит от способа вычисления лаг-члена (п.~\ref{sec:lag}); в частности, при вычислении лаг-члена по узлам сетки глобальный порядок оказывается равным $q$, а не $q-1$. Основные обозначения сведены в табл.~\ref{tab:notation}.

\begin{table}[ht]
\centering\small
\caption{Основные обозначения.}\label{tab:notation}
\begin{tabular}{ll}
\toprule
$\Phi_n(t)$ & лаг-член (интеграл памяти по $[0,t_n]$), формула~\eqref{eq:split}\\
$c,\ A,\ b$ & узлы, матрица и веса метода~\eqref{eq:method}\\
$g(t)=f(t,y(t))$, $g_k=g^{(k)}(t_n)$ & правая часть вдоль решения и её производные\\
$J=f_y$, $J'=\frac{\dd}{\dd t}f_y(t,y(t))$ & якобиан и его производная вдоль решения в точке $t_n$\\
$d_k$, $D_k$ & квадратурные дефекты стадий и выхода, формула~\eqref{eq:defects}\\
$\gamma_k=k!/\G(k+\alpha+1)$ & точные значения интегралов моментов\\
$B_m=b^{\top}A^{m-1}d_1$ & цепочечные дефекты\\
$C_1=b^{\top}d_2$, $C_2=b^{\top}\mathrm{diag}(c)d_1$ & коэффициенты при $Jg_2$ и $J'g_1$\\
$q$, $q(s,\alpha)$ & показатель локальной ошибки метода и его верхняя грань\\
$p$ & классический порядок (при $\alpha=1$: $q=p+1$)\\
\bottomrule
\end{tabular}
\end{table}

\begin{remark}[область применимости]\label{rem:scope}
Речь идёт о методах, в которых приращение на шаге строится только из стадий текущего шага; лаг-член может вычисляться любым способом. Схемы, использующие значения $f$ с прошлых шагов непосредственно в стадиях (аналоги двухшаговых методов Рунге--Кутты и общих линейных методов), теоремами этого раздела не покрываются, точно так же как классические барьеры Бутчера не ограничивают общие линейные методы.
\end{remark}

\subsection{Мотивирующий результат: несогласованность методов с переносом базовой точки}\label{sec:efork}

\paragraph{Конструкция EFORK.} Покажем исходное построение методов работы~\cite{Ghaffari2023}. Его авторы стремились перенести на дробный случай классическую схему вывода методов Рунге--Кутты: разложить точное и численное решения в ряды по шагу и приравнять коэффициенты. Роль формулы Тейлора при этом играет обобщённая формула Тейлора~\cite{OdibatShawagfeh2007}
\begin{equation}\label{eq:gentaylor}
y(t)=\sum_{k\ge0}\frac{(t-t_0)^{k\alpha}}{\G(k\alpha+1)}\,\bigl(D^{\alpha}\bigr)^ky(t_0),
\end{equation}
где $(D^{\alpha})^k$~--- $k$-кратное последовательное применение производной Капуто с началом в $t_0$, так что разложение ведётся по степеням $h^{\alpha}$, а не $h$. Метод записывается так же, как классический, с заменой $h$ на $h^{\alpha}$:
\[
K_i=h^{\alpha}f\Bigl(t_n+c_ih,\;y_n+\sum_{j<i}a_{ij}K_j\Bigr),\qquad y_{n+1}=y_n+\sum_i w_iK_i .
\]
Правая часть раскладывается смешанным рядом~--- по $(c_ih)^{\alpha}$ по времени и по Тейлору по $y$ (формула (8) в~\cite{Ghaffari2023}),~--- и сравнение коэффициентов при $h^{\alpha},h^{2\alpha},h^{3\alpha}$ даёт условия порядка. Для порядка $2\alpha$ (система (17) в~\cite{Ghaffari2023})
\[
\sum_iw_i=\frac{1}{\G(\alpha+1)},\qquad \sum_iw_ic_i^{\alpha}=\frac{\G(\alpha+1)}{\G(2\alpha+1)},
\]
например, двухстадийный метод $w_1=w_2=\frac{1}{2\G(\alpha+1)}$, $c_2=\bigl(2\G(\alpha+1)^2/\G(2\alpha+1)\bigr)^{1/\alpha}$, $a_{21}=c_2^{\alpha}/\G(\alpha+1)$; для порядка $3\alpha$ (система (22)) к ним добавляются условия $\sum_iw_ic_i^{2\alpha}=\G(2\alpha+1)/\G(3\alpha+1)$ и $\sum_iw_i(Ac^{\alpha})_i=\G(\alpha+1)/\G(3\alpha+1)$. Эта часть построения выполнена корректно: символьная проверка выявила лишь опечатку в промежуточных формулах (10)--(11), не влияющую на условия (22) (замечание~\ref{rem:typo}).

Формула~\eqref{eq:gentaylor} описывает решение вблизи начальной точки. Чтобы повторить шаг из $t_n$, авторы переносят начало производной Капуто в $t_n$, разбивая её на вклад прошлого и текущего шага:
\begin{equation}\label{eq:shift}
{}^{C}\!D^{\alpha}_{0}y(t)=M_n(t)+{}^{C}\!D^{\alpha}_{t_n}y(t),\qquad
M_n(t)=\frac{1}{\G(1-\alpha)}\int_0^{t_n}(t-s)^{-\alpha}y'(s)\dd s .
\end{equation}
Тогда на $[t_n,t_{n+1}]$ получается формально новая задача Коши с началом в $t_n$:
\[
{}^{C}\!D^{\alpha}_{t_n}y(t)=F_n(t,y):=f(t,y)-M_n(t),
\]
к которой применяется та же таблица Бутчера с $F_n$ вместо $f$; интеграл памяти $M_n$ приближается кусочно-линейной интерполяцией $y$ по узлам сетки. Заявлены глобальные порядки $2\alpha$ и $3\alpha$; численные эксперименты проведены на линейной задаче ${}^{C}\!D^{\alpha}y=-y+t^{4-\alpha}/\G(5-\alpha)$ с точным решением $y=t^4E_{\alpha,5}(-t^{\alpha})$ при $\alpha\in\{1/2,1/3,1/4\}$.

\paragraph{Скрытое допущение.} Построение опирается на то, что на каждом шаге решение $y(t_n+\tau)$ раскладывается по степеням $\tau^{\alpha},\tau^{2\alpha},\dots$ Это верно в начальной точке и, вообще говоря, неверно во внутренней. Рассмотрим простейший пример ${}^{C}\!D^{\alpha}y=\lambda y$, $y(0)=1$, с решением $y(t)=E_{\alpha}(\lambda t^{\alpha})=\sum_k\lambda^kt^{k\alpha}/\G(k\alpha+1)$. При $\alpha=\frac12$, $\lambda=-1$ оно выражается явно: $E_{1/2}(-\sqrt t)=e^{t}\operatorname{erfc}\sqrt t$, и
\[
y(t)=1-\frac{2}{\sqrt\pi}\,t^{1/2}+t-\frac{4}{3\sqrt\pi}\,t^{3/2}+\dots
\]
Вблизи $t=0$ решение раскладывается по дробным степеням, его производная $y'(t)\approx-1/\sqrt{\pi t}$ неограниченна, и естественной переменной служит именно $t^{\alpha}$~--- здесь разложение~\eqref{eq:gentaylor} законно. Но функция $E_{\alpha}$ целая, а $t^{\alpha}$ аналитична при $t>0$, поэтому в любой точке $t_n>0$ решение аналитично и раскладывается по \emph{целым} степеням $\tau=t-t_n$: $y(t_n+\tau)=y(t_n)+y'(t_n)\tau+\frac12y''(t_n)\tau^2+\dots$ Записать его рядом по $\tau^{1/2}$ нельзя: любой ненулевой член с $\tau^{1/2}$ сделал бы функцию недифференцируемой в $\tau=0$. Негладкость решения сосредоточена в начальной точке, и разложение~\eqref{eq:gentaylor} переносится на внутренний шаг неправомерно.

Применяя к целому ряду производную Капуто с началом в $t_n$, получаем
\begin{equation}\label{eq:caputo-interior}
{}^{C}\!D^{\alpha}_{t_n}y(t_n+\tau)=\frac{y'(t_n)}{\G(2-\alpha)}\,\tau^{1-\alpha}+\frac{y''(t_n)}{\G(3-\alpha)}\,\tau^{2-\alpha}+\dots,
\end{equation}
т.\,е. степени $1-\alpha,\,2-\alpha,\dots$, тогда как метод подгоняет коэффициенты при степенях $\alpha,2\alpha,3\alpha,\dots$ Эти множества показателей совпадают лишь при $1-\alpha=k\alpha$, т.\,е. при $\alpha=1/(k+1)$; при прочих $\alpha$ метод добросовестно подгоняет не тот ряд. Отметим также следующий тонкий момент. Левая часть~\eqref{eq:shift}, равная $f(t,y(t))$, во внутренней точке гладкая, а слагаемое~\eqref{eq:caputo-interior} содержит член $+y'(t_n)\tau^{1-\alpha}/\G(2-\alpha)$; следовательно, интеграл памяти $M_n(t_n+\tau)$ содержит тот же член с обратным знаком, и в сумме они сокращаются. Перенос базовой точки разрезает гладкую функцию на две негладкие части, а метод аппроксимирует лишь одну из них, обращаясь с $F_n=f-M_n$ как с гладкой правой частью новой задачи. В частности, при точной памяти $F_n(t_n,y(t_n))=0$, и первая стадия $K_1$ вообще не даёт вклада.

\begin{theorem}[о необходимом условии согласованности EFORK]\label{thm:efork}
Пусть $n\ge1$, $y\in C^2[t_n,t_{n+1}]$, $f$ липшицева по второму аргументу, и интеграл памяти $M_n$ вычислен точно. Тогда для любого явного метода EFORK
\[
y_{n+1}-y(t_{n+1})=h\,y'(t_n)\bigl(R(\alpha)-1\bigr)+O(h^{1+\alpha}),\qquad
R(\alpha)=\frac{1}{\G(2-\alpha)}\sum_{i=1}^{s}w_ic_i^{\,1-\alpha}.
\]
\end{theorem}
\begin{proof}
По~\eqref{eq:caputo-interior} $F_n(t_n+\tau,y(t_n+\tau))=\frac{y'(t_n)}{\G(2-\alpha)}\tau^{1-\alpha}+O(\tau^{2-\alpha})$. По индукции по $i$ получаем $K_j=O(h)$, откуда $z_i:=y_n+\sum_j a_{ij}K_j=y(t_n+c_ih)+O(h)$, и по липшицевости
\[
K_i=h^{\alpha}\Bigl[\frac{y'(t_n)}{\G(2-\alpha)}(c_ih)^{1-\alpha}+O(h^{2-\alpha})+O(h)\Bigr]
=\frac{h\,y'(t_n)\,c_i^{1-\alpha}}{\G(2-\alpha)}+O(h^{1+\alpha}).
\]
Суммируя с весами $w_i$, находим $y_{n+1}-y_n=h\,y'(t_n)R(\alpha)+O(h^{1+\alpha})$, тогда как точное приращение равно $y(t_{n+1})-y(t_n)=h\,y'(t_n)+O(h^2)$.
\end{proof}

В опубликованной реализации интеграл памяти вычисляется не точно, а по кусочно-линейной интерполяции решения. Покажем, что это не меняет главного члена.

\begin{remark}[приближённый интеграл памяти]\label{rem:memory}
Пусть $t_n\ge t_*>0$, значения $y(t_m)$, $m\le n$, точные (как и положено при определении локальной ошибки), а решение обладает типичной для уравнений с производной Капуто регулярностью: $y\in C[0,T]\cap C^2(0,T]$, $|y(s)-y(0)|\le Cs^{\alpha}$, $|y''(s)|\le Cs^{\alpha-2}$~\cite{Brunner2004,Diethelm2010}. Обозначим через $\tilde y$ кусочно-линейный интерполянт $y$ по узлам сетки, а через $\tilde M_n$~--- интеграл~\eqref{eq:shift} с $\tilde y'$ вместо $y'$. Поскольку $\tilde y-y$ обращается в нуль в узлах, интегрирование по частям на каждом интервале даёт
\[
\tilde M_n(t)-M_n(t)=-\frac{\alpha}{\G(1-\alpha)}\int_0^{t_n}(t-s)^{-\alpha-1}\bigl(\tilde y(s)-y(s)\bigr)\dd s,\qquad t\in[t_n,t_{n+1}].
\]
На первом интервале $|\tilde y-y|\le Ch^{\alpha}$, на интервале $[t_k,t_{k+1}]$, $k\ge1$,~--- $|\tilde y-y|\le Ch^2(kh)^{\alpha-2}$, а на последнем интервале $|\tilde y(s)-y(s)|\le Ch\,(t_n-s)$. При $s\le t_n/2$ ядро ограничено, и вклад этой части не превосходит $Ch\bigl(h^{\alpha}+h^{\alpha}\sum_kk^{\alpha-2}\bigr)=O(h^{1+\alpha})$; на $[t_n/2,t_{n-1}]$ имеем $t-s\ge h$ и вклад $O\bigl(h^2\int_h^{\infty}u^{-\alpha-1}\dd u\bigr)=O(h^{2-\alpha})$; на последнем интервале вклад $O\bigl(h\int_0^h u^{-\alpha}\dd u\bigr)=O(h^{2-\alpha})$. Итак, $|\tilde M_n-M_n|=O(h^{\min(1+\alpha,\,2-\alpha)})$ равномерно на $[t_n,t_{n+1}]$, и каждая стадия получает добавку $h^{\alpha}O(h^{\min(1+\alpha,2-\alpha)})=O(h^{\min(1+2\alpha,\,2)})$, что входит в остаток $O(h^{1+\alpha})$. Следовательно, утверждение теоремы~\ref{thm:efork} справедливо и для опубликованной реализации методов EFORK. Оценка подтверждается численно: для $y=1+t^{\alpha}+t^2$, $t_n=0.5$ наблюдаемый порядок погрешности $\tilde M_n$ равен $1.27\to1.29$ при $\alpha=0.3$ (теория $1.3$) и $1.03\to1.22$ при $\alpha=0.7$ (теория $1.3$, выход на асимптотику медленнее).
\end{remark}

\begin{corollary}
Необходимое условие согласованности EFORK на внутренних шагах: $\sum_i w_ic_i^{1-\alpha}=\G(2-\alpha)$. Если $1-\alpha=k\alpha$ при натуральном $k\le p-1$, то правая часть условия порядка $\sum_iw_ic_i^{k\alpha}=\G(k\alpha+1)/\G((k+1)\alpha+1)$ из~\cite{Ghaffari2023} равна $\G(2-\alpha)/\G(2)=\G(2-\alpha)$, и оно совпадает с условием согласованности. Поэтому при $\alpha\in\{1/2,1/3,\dots,1/p\}$ метод согласован случайно~--- и именно такие значения использованы в экспериментах~\cite{Ghaffari2023}. При прочих $\alpha$ условие является дополнительным: условия порядка EFORK сами по себе согласованности не обеспечивают. Это не исключает, что коэффициенты отдельного метода удовлетворят ему случайно или будут подобраны специально; утверждать несогласованность можно лишь для конкретных методов. Для опубликованных методов~\cite{Ghaffari2023} величина $R(\alpha)$ вычисляется явно (рис.~\ref{fig:efork}а, табл.~\ref{tab:efork}): двухстадийный метод согласован лишь при $\alpha=1/2$ (уже при $\alpha=1/4$ имеем $R\approx3.8$), трёхстадийный с $c=(0,\frac12,1)$~--- лишь при $\alpha\in\{1/3,1/2\}$.
\end{corollary}

\begin{remark}\label{rem:typo}
В формулах (10)--(11) работы~\cite{Ghaffari2023} коэффициенты при $f^2f_{yy}$ и $ff^{\alpha,1}_{t,y}$ должны равняться $\G(2\alpha+1)/(2\G(\alpha+1)^2)$ и $\G(2\alpha+1)/\G(\alpha+1)^2$ (а не $1/2$ и $1$); промежуточная формула на с.~6 там записана верно, и условия порядка (22) от этой опечатки не зависят.
\end{remark}

Таким образом, правильный подход к построению методов состоит не в переносе базовой точки, а в разбиении~\eqref{eq:split}: на текущем шаге интегрируется гладкая функция $g(t)=f(t,y(t))$, а вся дробность сосредоточена в ядре квадратуры.

\subsection{Структура локальной ошибки}

\paragraph{Точное решение на шаге.} Обозначим $g(t)=f(t,y(t))$, $g_k=g^{(k)}(t_n)$. Подставляя в~\eqref{eq:split} ряд Тейлора $g(t_n+\theta h)=\sum_kg_kh^k\theta^k/k!$ и вычисляя интеграл от монома через бета-функцию,
\begin{equation}\label{eq:beta}
\frac{1}{\G(\alpha)}\int_0^{c}(c-\theta)^{\alpha-1}\theta^k\dd\theta=\frac{k!\,c^{k+\alpha}}{\G(k+\alpha+1)},
\end{equation}
получаем
\[
y(t_n+ch)=\Phi_n(t_n+ch)+h^{\alpha}\sum_{k\ge0}g_kh^k\,\frac{c^{k+\alpha}}{\G(k+\alpha+1)} .
\]

\paragraph{Квадратурные дефекты.} Если бы стадии вычислялись по точным значениям $F_j=g(t_n+c_jh)$, стадия $i$ дала бы $\Phi_n(t_n+c_ih)+h^{\alpha}\sum_kg_kh^k(Ac^k)_i/k!$. Разность с точным значением определяет \emph{квадратурные дефекты} стадий (векторы) и выхода (числа):
\begin{equation}\label{eq:defects}
d_{k}=Ac^{k}-\frac{k!}{\G(k+\alpha+1)}\,c^{\,k+\alpha},\qquad
D_k=b^{\top}c^{k}-\frac{k!}{\G(k+\alpha+1)} .
\end{equation}
Величина $(d_k)_i$ показывает, насколько квадратура стадии $i$ с весами $a_{ij}$ ошибается на мономе $\theta^k$, проинтегрированном с ядром $(c_i-\theta)^{\alpha-1}$; $D_k$~--- то же для итоговой квадратуры с весами $b$. По условию~\eqref{eq:rowsum} $d_0=0$, а первая компонента любого $d_k$ равна нулю, поскольку $c_1=0$ и первая строка $A$ нулевая. Положим также
\[
\gamma_0=\frac{1}{\G(\alpha+1)},\qquad\gamma_1=\frac{1}{\G(\alpha+2)},\qquad\gamma_2=\frac{2}{\G(\alpha+3)}
\]
(точные значения~\eqref{eq:beta} при $c=1$, $k=0,1,2$), так что $D_k=b^{\top}c^k-\gamma_k$.

\paragraph{Ошибки стадий.} На самом деле $F_j$ вычисляется в $Y_j$, а не в точке решения. Обозначим $e_i=Y_i-y(t_n+c_ih)$, $J=f_y(t_n,y(t_n))$, $J'=\frac{\dd}{\dd t}f_y(t,y(t))\big|_{t_n}$. Поскольку якобиан берётся в точке стадии, $f_y(t_n+c_jh,y(t_n+c_jh))=J+hc_jJ'+\dots$, и
\[
F_j=g(t_n+c_jh)+\bigl(J+hc_jJ'+\dots\bigr)e_j+\tfrac12f_{yy}(e_j,e_j)+\dots
\]
Для вектора ошибок стадий отсюда следует рекуррентное соотношение (для систем произведения $AJ$ понимаются как кронекеровы: $A$ действует на номера стадий, $J$~--- на компоненты решения)
\begin{equation}\label{eq:stageerr}
e=h^{\alpha}\sum_{k\ge1}\frac{h^k g_k}{k!}\,d_k+h^{\alpha}AJe+h^{1+\alpha}A\,\mathrm{diag}(c)\,J'e+h^{\alpha}A\Bigl[\tfrac12f_{yy}(e,e)+O\bigl(h^2|e|\bigr)\Bigr],
\end{equation}
а для локальной ошибки~--- то же выражение с заменой $A$ на $b^{\top}$ и $d_k$ на $D_k$ (и с членом $k=0$). Итерируя~\eqref{eq:stageerr}, последовательно получаем:
\begin{itemize}
\item в нулевом приближении $e=h^{1+\alpha}g_1d_1+O(h^{2+\alpha})$: член $k=0$ отсутствует, так как $d_0=0$, и ошибки стадий начинаются с $h^{1+\alpha}$;
\item вклад выхода без учёта ошибок стадий: $h^{\alpha}D_0g_0+h^{1+\alpha}D_1g_1+\frac12h^{2+\alpha}D_2g_2+\dots$;
\item ошибка стадий, прошедшая через якобиан: $h^{\alpha}b^{\top}J\,h^{1+\alpha}g_1d_1=h^{1+2\alpha}\,(b^{\top}d_1)\,Jg_1$;
\item каждое следующее прохождение через $h^{\alpha}AJ$ добавляет множитель $h^{\alpha}$: возникает цепочка $h^{1+(m+1)\alpha}\,b^{\top}A^{m-1}d_1\,J^mg_1$;
\item член $\frac12h^{2+\alpha}g_2d_2$ ошибки стадий, прошедший через $J$, даёт $\frac12h^{2+2\alpha}\,(b^{\top}d_2)\,Jg_2$; поправка $hc_jJ'$ к якобиану, умноженная на $h^{1+\alpha}g_1(d_1)_j$, даёт $h^{2+2\alpha}\,(b^{\top}\mathrm{diag}(c)d_1)\,J'g_1$;
\item нелинейные члены начинаются с $h^{\alpha}f_{yy}(e,e)\sim h^{\alpha}(h^{1+\alpha})^2=h^{2+3\alpha}$.
\end{itemize}

\begin{proposition}[структура локальной ошибки]\label{prop:structure}
Для метода~\eqref{eq:method}, \eqref{eq:rowsum} локальная ошибка имеет вид
\begin{align*}
\delta_n={}&h^{\alpha}D_0\,g_0+h^{1+\alpha}D_1\,g_1+\tfrac12h^{2+\alpha}D_2\,g_2
+\sum_{m=1}^{s-1}h^{1+(m+1)\alpha}B_m\,J^{m}g_1\\
&+h^{2+2\alpha}\bigl(\tfrac12C_1\,Jg_2+C_2\,J'g_1\bigr)+\rho_n,
\qquad B_m=b^{\top}A^{m-1}d_1,\quad C_1=b^{\top}d_2,\quad C_2=b^{\top}\mathrm{diag}(c)\,d_1,
\end{align*}
где $\rho_n$ состоит из членов вида $h^{k+m\alpha}\times$(элементарный дифференциал) с $k\ge3$, либо с $k=2$ и $m\ge3$; в частности, если $f\in C^3$ в окрестности решения и $y\in C^3[t_n,t_n+\hat ch]$, то $|\rho_n|\le K\,h^{\min(3+\alpha,\,2+3\alpha)}$ с постоянной $K$, зависящей от норм производных $f$ до третьего порядка, от $\max_i|b_i|,\max_{ij}|a_{ij}|$ и от $\hat c$, но не от $h$. Если $y\in C^3[t_*,T+\hat ch_0]$ при некотором $t_*>0$, то постоянная $K$ равномерна по всем шагам с $t_n\ge t_*$ и $h\le h_0$.
\end{proposition}

\begin{proof}
\emph{Шаг 1: априорная оценка.} Поскольку $f\in C^3$, имеем $F_j=g(t_n+c_jh)+M_je_j$, где $M_j=\int_0^1f_y(t_n+c_jh,\,y(t_n+c_jh)+\tau e_j)\dd\tau$. Вычитая из~\eqref{eq:method} точное соотношение~\eqref{eq:split} и применяя формулу Тейлора с остаточным членом к $g\in C^3$ вместе с~\eqref{eq:beta}, получаем
\[
e_i=h^{\alpha}\Bigl(h g_1(d_1)_i+\tfrac{h^2}{2}g_2(d_2)_i\Bigr)+h^{\alpha}\sum_{j<i}a_{ij}M_je_j+r_i,\qquad |r_i|\le K_1h^{3+\alpha}.
\]
Метод явный, поэтому $e_1=0$, а $e_i$ выражается через $e_1,\dots,e_{i-1}$. Индукцией по $i$ получаем $|e_i|\le K_2h^{1+\alpha}$ при $h\le h_0$.
\emph{Шаг 2: разложение $M_j$.} По формуле Тейлора $M_j=J+hc_jJ'+R_j$, где $|R_j|\le K_3(h^2+|e_j|)\le K_4h^{1+\alpha}$. Следовательно,
\[
e=h^{1+\alpha}g_1d_1+\tfrac12h^{2+\alpha}g_2d_2+h^{\alpha}AJe+h^{1+\alpha}A\,\mathrm{diag}(c)J'e+\varepsilon,\qquad |\varepsilon|\le K_5h^{\min(3+\alpha,\,2+3\alpha)},
\]
так как $h^{\alpha}R_je_j=O(h^{2+3\alpha})$.
\emph{Шаг 3: конечная итерация.} Матрица $A$ нильпотентна, поэтому $I-h^{\alpha}AJ$ обратима и $(I-h^{\alpha}AJ)^{-1}=\sum_{m=0}^{s-1}h^{m\alpha}(AJ)^m$~--- конечная сумма. Подставляя первое приближение $e^{(0)}=h^{1+\alpha}g_1d_1$ в член с $J'$, находим
\[
e=\sum_{m=0}^{s-1}h^{m\alpha}(AJ)^m\Bigl[h^{1+\alpha}g_1d_1+\tfrac12h^{2+\alpha}g_2d_2+h^{2+2\alpha}A\,\mathrm{diag}(c)\,J'g_1d_1\Bigr]+O\bigl(h^{\min(3+\alpha,\,2+3\alpha)}\bigr).
\]
\emph{Шаг 4: локальная ошибка.} Аналогично
\[
\delta_n=h^{\alpha}\sum_{k=1}^{2}\frac{h^kg_k}{k!}D_k+h^{\alpha}b^{\top}Je+h^{1+\alpha}b^{\top}\mathrm{diag}(c)J'e+O\bigl(h^{\min(3+\alpha,\,2+3\alpha)}\bigr).
\]
Подставляя $e$, получаем: при $m\ge0$ из первого слагаемого~--- члены $h^{1+(m+2)\alpha}b^{\top}A^md_1J^{m+1}g_1$, т.\,е. цепочку $B_{m+1}$; из второго~--- $\frac12h^{2+2\alpha}C_1Jg_2$ при $m=0$ и члены порядка не ниже $h^{2+3\alpha}$ при $m\ge1$; из третьего и из слагаемого $h^{1+\alpha}b^{\top}\mathrm{diag}(c)J'e$ с $e^{(0)}$~--- член $h^{2+2\alpha}C_2J'g_1$ и члены порядка не ниже $h^{2+3\alpha}$. Все прочие слагаемые входят в $\rho_n$.
\end{proof}

Явный вид предложения дополнительно проверен символьно-численным генератором условий порядка (п.~\ref{sec:generator}). Для скалярных уравнений и для систем условия совпадают по крайней мере до показателя $2+3\alpha$; так как $2+3\alpha>\max(2+2\alpha,1+4\alpha)$ при $\alpha<1$, все результаты этого раздела при $s\le4$ справедливы и для систем. Сводка младших членов приведена в табл.~\ref{tab:terms}.

\begin{table}[ht]
\centering\small
\caption{Младшие члены локальной ошибки: коэффициент, элементарный дифференциал, степень $h$ и условие обнуления.}\label{tab:terms}
\begin{tabular}{llll}
\toprule
член & дифференциал & степень $h$ & условие обнуления\\
\midrule
$D_0$ & $g_0=f$ & $\alpha$ & $\sum_ib_i=\gamma_0$\\
$D_1$ & $g_1$ & $1+\alpha$ & $\sum_ib_ic_i=\gamma_1$\\
$D_2$ & $g_2$ & $2+\alpha$ & $\sum_ib_ic_i^2=\gamma_2$\\
$B_m$ & $J^mg_1$ & $1+(m+1)\alpha$ & $b^{\top}A^{m-1}d_1=0$\\
$C_1$ & $Jg_2$ & $2+2\alpha$ & $\sum_ib_i(d_2)_i=0$\\
$C_2$ & $J'g_1$ & $2+2\alpha$ & $\sum_ib_ic_i(d_1)_i=0$\\
\bottomrule
\end{tabular}
\end{table}

Показатель метода равен степени первого члена, который не удалось обнулить. Например, при $\alpha=\frac12$ член $D_0$ стоит при $h^{0.5}$, $D_1$~--- при $h^{1.5}$, $B_1$~--- при $h^{2}$, $D_2$ и $B_2$~--- при $h^{2.5}$, $B_3$, $C_1$ и $C_2$~--- при $h^{3}$. При $\alpha=1$ все показатели становятся целыми и <<слипаются>>: например, $D_1$ и $B_1$ оба стоят при $h^2$ и вместе дают одно классическое условие порядка. При $\alpha<1$ они расходятся по разным степеням $h$, и каждый приходится обнулять отдельно; в этом и состоит источник дробных барьеров.

\subsection{Барьеры}

Перейдём к основным результатам. Начнём с наблюдения, которое служит главным рычагом всех доказательств.

\begin{lemma}[о полной цепочке]\label{lem:chain}
Для явного $s$-стадийного метода класса~\eqref{eq:rowsum} $B_m=0$ при $m\ge s$, а последнее звено цепочки равно
\[
B_{s-1}=b_s\,a_{s,s-1}a_{s-1,s-2}\cdots a_{32}\,(d_1)_2,\qquad (d_1)_2=-\frac{c_2^{1+\alpha}}{\G(\alpha+2)} .
\]
\end{lemma}
\begin{proof}
Матрица $A$ строго нижнетреугольная, поэтому $A^s=0$, а у $A^{s-1}$ ненулевым может быть лишь элемент $(A^{s-1})_{s1}=a_{s,s-1}\cdots a_{21}$. Значит, $A^{s-1}d_1$ <<читает>> только первую компоненту $d_1$, равную нулю, и $B_s=0$. Аналогично у $A^{s-2}$ ненулевыми могут быть лишь элементы $(s,1)$, $(s-1,1)$ и $(s,2)$, откуда $B_{s-1}=b^{\top}A^{s-2}d_1=b_s(A^{s-2})_{s2}(d_1)_2$ с $(A^{s-2})_{s2}=a_{s,s-1}\cdots a_{32}$. Наконец, $(d_1)_2=a_{21}c_1-c_2^{1+\alpha}/\G(\alpha+2)$ и $c_1=0$.
\end{proof}

Формула для $(d_1)_2$ имеет простой смысл: вторая стадия <<видит>> правую часть лишь в одной точке $\theta=0$, где расположена первая стадия, и потому не может проинтегрировать даже линейную функцию $\theta$. Это фундаментальное свойство явных методов, из которого вырастает вся цепочка $B_m$. Обнулить $B_{s-1}$ можно, лишь положив $b_s=0$ или $c_2=0$ (что фактически уменьшает число стадий), либо <<разорвав цепочку>>, т.\,е. положив $a_{i+1,i}=0$ для некоторого $i$.

Второй инструмент~--- классическое утверждение о нулях обобщённых многочленов.

\begin{lemma}[правило знаков Декарта--Лагерра~\cite{PolyaSzego}]\label{lem:descartes}
Пусть $p_1<p_2<\dots<p_r$~--- вещественные числа, $\mu_1,\dots,\mu_r$~--- вещественные коэффициенты, не все равные нулю. Тогда число нулей функции $\sum_{i}\mu_ix^{p_i}$ на интервале $(0,\infty)$, подсчитанных с учётом кратности, не превосходит числа перемен знака в последовательности $\mu_1,\dots,\mu_r$, из которой предварительно удалены нулевые члены.
\end{lemma}

Для обычных многочленов это правило Декарта: например, у $x^2-3x+2$ последовательность знаков $(+,-,+)$ содержит две перемены, и положительных корней два ($1$ и $2$); Лагерр распространил его на произвольные вещественные показатели.

\begin{lemma}[о мерах на узлах]\label{lem:measures}
Не существует набора точек $0<x_1<\dots<x_k$ и весов $\rho_1,\dots,\rho_k\in\mathbb R$ таких, что
\begin{enumerate}
\item[(M2)] $k\le2$ и $\sum\rho_i=\gamma_1$, $\sum\rho_ix_i=\gamma_2$, $\sum\rho_ix_i^{\alpha}=0$;
\item[(M3)] $k\le3$ и выполнены условия (M2), а также $\sum\rho_ix_i^{1+\alpha}=0$.
\end{enumerate}
\end{lemma}
\begin{proof}
Рассмотрим (M3). Пространство функций $\psi(x)=\mu_0+\mu_2x^{\alpha}+\mu_1x+\mu_3x^{1+\alpha}$ четырёхмерно, поэтому найдётся $\psi\not\equiv0$, обращающаяся в нуль во всех $x_i$ (коэффициентов четыре, условий не более трёх). Умножая $\psi(x_i)=0$ на $\rho_i$ и суммируя, получаем $0=\mu_0\gamma_1+\mu_2\cdot0+\mu_1\gamma_2+\mu_3\cdot0$, т.\,е. $\mu_0=-\frac{\gamma_2}{\gamma_1}\mu_1=-\frac{2}{\alpha+2}\mu_1$. Если $\mu_1\neq0$, то в упорядоченной по показателям $0<\alpha<1<1+\alpha$ последовательности $(\mu_0,\mu_2,\mu_1,\mu_3)$ первый и третий члены имеют противоположные знаки; значит, на участке $(\mu_0,\mu_2,\mu_1)$ происходит ровно одна перемена знака, и всего перемен не больше двух. По лемме~\ref{lem:descartes} $\psi$ имеет не более двух положительных нулей~--- противоречие при $k=3$; при $k<3$ можно добавить узлы с нулевыми весами. Если $\mu_1=0$, то $\mu_0=0$ и $\psi=x^{\alpha}(\mu_2+\mu_3x)$ имеет не более одного положительного нуля. Случай (M2) разбирается так же с $\mu_3=0$: перемена знака одна, а нулей требуется два.
\end{proof}

\paragraph{Две стадии.} Метод задаётся таблицей $c=(0,c_2)$, $a_{21}=c_2^{\alpha}/\G(\alpha+1)$, $b=(b_1,b_2)$; свободны $c_2,b_1,b_2$. Компоненты дефекта: $(d_1)_1=0$, $(d_1)_2=-c_2^{1+\alpha}/\G(\alpha+2)$.

\begin{theorem}[барьер для двух стадий]\label{thm:s2}
$q(2,\alpha)=1+2\alpha$.
\end{theorem}
\begin{proof}
По табл.~\ref{tab:terms} ниже показателя $1+2\alpha$ стоят только $D_0$ и $D_1$; они обнуляются выбором $b_2c_2=\gamma_1$, $b_1=\gamma_0-b_2$ при любом $c_2>0$. Член при $h^{1+2\alpha}$ имеет коэффициент $B_1=b_2(d_1)_2=-b_2c_2\cdot c_2^{\alpha}/\G(\alpha+2)=-\gamma_1c_2^{\alpha}/\G(\alpha+2)\neq0$, так как $c_2=0$ противоречит условию $D_1$.
\end{proof}

При $\alpha=1$ получаем $q=3$, т.\,е. $p=2$~--- классический двухстадийный метод второго порядка.

\paragraph{Три стадии.} Метод задаётся таблицей
\[
c=(0,c_2,c_3),\qquad A=\begin{pmatrix}0&0&0\\a_{21}&0&0\\a_{31}&a_{32}&0\end{pmatrix},\qquad b=(b_1,b_2,b_3),
\]
где $a_{21}=c_2^{\alpha}/\G(\alpha+1)$, $a_{31}+a_{32}=c_3^{\alpha}/\G(\alpha+1)$; свободны шесть параметров $c_2,c_3,a_{32},b_1,b_2,b_3$. Компоненты дефектов и произведения $Ad_1$:
\begin{gather*}
(d_1)_1=0,\qquad (d_1)_2=-\frac{c_2^{1+\alpha}}{\G(\alpha+2)},\qquad (d_1)_3=a_{32}c_2-\frac{c_3^{1+\alpha}}{\G(\alpha+2)},\\
(Ad_1)_1=(Ad_1)_2=0,\qquad (Ad_1)_3=a_{32}(d_1)_2 .
\end{gather*}
Условия, которые предстоит обнулить, сведены в табл.~\ref{tab:s3cond}. Неизвестных шесть, условий пять, так что наивный подсчёт обещает решение; мешает, однако, не число уравнений, а их структура.

\begin{table}[ht]
\centering\small
\caption{Трёхстадийный метод: условия и степени $h$; справа~--- значения показателей при $\alpha=0.3$ и $\alpha=0.7$.}\label{tab:s3cond}
\begin{tabular}{llccc}
\toprule
условие & формула & степень $h$ & $\alpha=0.3$ & $\alpha=0.7$\\
\midrule
$D_0$ & $b_1+b_2+b_3=\gamma_0$ & $\alpha$ & 0.3 & 0.7\\
$D_1$ & $b_2c_2+b_3c_3=\gamma_1$ & $1+\alpha$ & 1.3 & 1.7\\
$B_1$ & $b_2(d_1)_2+b_3(d_1)_3=0$ & $1+2\alpha$ & 1.6 & 2.4\\
$B_2$ & $b_3a_{32}(d_1)_2=0$ & $1+3\alpha$ & 1.9 & 3.1\\
$D_2$ & $b_2c_2^2+b_3c_3^2=\gamma_2$ & $2+\alpha$ & 2.3 & 2.7\\
\bottomrule
\end{tabular}
\end{table}

\begin{theorem}[барьер для трёх стадий]\label{thm:s3}
$q(3,\alpha)=\max(2+\alpha,\,1+3\alpha)$.
\end{theorem}
\begin{proof}
\emph{Оценка сверху.} Превзойти указанное значение можно, лишь обнулив одновременно все пять условий табл.~\ref{tab:s3cond}. По лемме~\ref{lem:chain} $B_2=b_3a_{32}(d_1)_2$, и равенство $B_2=0$ требует обращения в нуль одного из множителей. Если $c_2=0$, то $Y_2=Y_1$ и метод фактически двухстадийный; если $b_3=0$~--- тоже. В обоих случаях $B_1\neq0$ по теореме~\ref{thm:s2}. Остаётся случай $a_{32}=0$. Если при этом $c_3=0$, то по~\eqref{eq:rowsum} и $a_{31}=0$, т.\,е. $Y_3=Y_1$, и метод снова фактически двухстадийный; поэтому далее $c_3>0$. Тогда $(d_1)_3=-c_3^{1+\alpha}/\G(\alpha+2)$, и условие $B_1$ после умножения на $-\G(\alpha+2)$ принимает вид $b_2c_2^{1+\alpha}+b_3c_3^{1+\alpha}=0$. Обозначив $\rho_i=b_ic_i$, получаем систему
\[
\text{(I)}\ \rho_2+\rho_3=\gamma_1,\qquad \text{(II)}\ \rho_2c_2+\rho_3c_3=\gamma_2,\qquad \text{(III)}\ \rho_2c_2^{\alpha}+\rho_3c_3^{\alpha}=0,
\]
т.\,е. систему (M2) леммы~\ref{lem:measures}. Покажем её несовместность и элементарно. Если $c_2=c_3$, то (III) даёт $\rho_2+\rho_3=0$ вопреки (I). Пусть $r=c_2/c_3\neq1$. Из (III) $\rho_3=-\rho_2r^{\alpha}$; подставляя в (I) и (II), находим $\rho_2(1-r^{\alpha})=\gamma_1$ и $\rho_2c_3(r-r^{\alpha})=\gamma_2$, откуда
\[
c_3\,\frac{r-r^{\alpha}}{1-r^{\alpha}}=\frac{\gamma_2}{\gamma_1}=\frac{2}{\alpha+2}>0 .
\]
Но при $0<\alpha<1$ и $r<1$ имеем $r^{\alpha}>r$, так что числитель отрицателен, а знаменатель положителен; при $r>1$~--- наоборот. Левая часть всегда отрицательна, противоречие.

\emph{Достижимость.} Обнулить одновременно $D_2$ и $B_2$ нельзя, но можно обнулить любой один из них.
Семейство A (<<разрыв цепочки>>): $a_{32}=0$, $b_3=-b_2(c_2/c_3)^{1+\alpha}$, $b_2=\gamma_1/\bigl(c_2-c_2^{1+\alpha}c_3^{-\alpha}\bigr)$, $b_1=\gamma_0-b_2-b_3$. Тогда $D_0=D_1=B_1=B_2=0$, член $B_3$ отсутствует, так как $A^2d_1=0$, и первым ненулевым остаётся $D_2$: $q_A=2+\alpha$.
Семейство B (<<цепочка>>): $b$ из условий $D_0=D_1=D_2=0$ и $a_{32}=\bigl(b_2c_2^{1+\alpha}+b_3c_3^{1+\alpha}\bigr)/\bigl(\G(\alpha+2)\,b_3c_2\bigr)$, что даёт $B_1=0$; первым ненулевым остаётся $B_2$: $q_B=1+3\alpha$.
Выбирая лучшее из двух семейств, получаем $q(3,\alpha)=\max(q_A,q_B)$.
\end{proof}

Правило выбора наглядно видно из табл.~\ref{tab:s3cond}. Поскольку обнулить можно лишь один из двух членов $D_2$, $B_2$, следует гасить тот, что стоит ниже, и показатель равен степени верхнего. При $\alpha=0.3$ член $B_2$ (показатель $1.9$) лежит ниже $D_2$ ($2.3$); семейство A гасит $B_2$ и даёт $q=2.3$, семейство B оставляет $B_2$ и даёт лишь $1.9$. При $\alpha=0.7$ порядок обратный: $D_2$ ($2.7$) ниже $B_2$ ($3.1$), и выигрывает семейство B с $q=3.1$. Ветви сравниваются при $2+\alpha=1+3\alpha$, т.\,е. при $\alpha=\frac12$, где $q=2.5$; при $\alpha=1$ получаем $q=4$, т.\,е. $p=3$~--- классический трёхстадийный метод.

\paragraph{Четыре стадии.} Метод задаётся таблицей
\[
c=(0,c_2,c_3,c_4),\qquad A=\begin{pmatrix}0&0&0&0\\a_{21}&0&0&0\\a_{31}&a_{32}&0&0\\a_{41}&a_{42}&a_{43}&0\end{pmatrix},\qquad b=(b_1,b_2,b_3,b_4),
\]
с суммами строк $a_{21}=c_2^{\alpha}/\G(\alpha+1)$, $a_{31}+a_{32}=c_3^{\alpha}/\G(\alpha+1)$, $a_{41}+a_{42}+a_{43}=c_4^{\alpha}/\G(\alpha+1)$; свободны десять параметров $c_2,c_3,c_4,a_{32},a_{42},a_{43},b_1,\dots,b_4$ (элементы первого столбца определяются суммами строк и в дефекты не входят, поскольку умножаются на $c_1=0$). Компоненты дефектов:
\begin{gather*}
(d_1)_2=-\tfrac{c_2^{1+\alpha}}{\G(\alpha+2)},\quad (d_1)_3=a_{32}c_2-\tfrac{c_3^{1+\alpha}}{\G(\alpha+2)},\quad (d_1)_4=a_{42}c_2+a_{43}c_3-\tfrac{c_4^{1+\alpha}}{\G(\alpha+2)},\\
(d_2)_2=-\tfrac{2c_2^{2+\alpha}}{\G(\alpha+3)},\quad (d_2)_3=a_{32}c_2^2-\tfrac{2c_3^{2+\alpha}}{\G(\alpha+3)},\quad (d_2)_4=a_{42}c_2^2+a_{43}c_3^2-\tfrac{2c_4^{2+\alpha}}{\G(\alpha+3)},
\end{gather*}
а цепочки равны
\[
B_2=b^{\top}Ad_1=b_3a_{32}(d_1)_2+b_4\bigl[a_{42}(d_1)_2+a_{43}(d_1)_3\bigr],\qquad B_3=b^{\top}A^2d_1=b_4a_{43}a_{32}(d_1)_2 .
\]
При четырёх стадиях впервые существенны члены при $h^{2+2\alpha}$. Их два, и происхождение у них разное. Член $C_1=\sum_ib_i(d_2)_i$ возникает из того, что стадии неточно интегрируют квадратичную часть $g$: слагаемое $h^{2+\alpha}\frac{g_2}{2}(d_2)_i$ ошибки стадии, пройдя через якобиан, даёт $h^{\alpha}\sum_ib_iJh^{2+\alpha}\frac{g_2}{2}(d_2)_i=\frac12h^{2+2\alpha}C_1Jg_2$. Член $C_2=\sum_ib_ic_i(d_1)_i$ связан с изменением якобиана вдоль шага: поправка $hc_iJ'$, умноженная на ошибку стадии $h^{1+\alpha}g_1(d_1)_i$, даёт $h^{\alpha}\sum_ib_i\,hc_iJ'\,h^{1+\alpha}g_1(d_1)_i=h^{2+2\alpha}C_2J'g_1$. Для линейной задачи с постоянными коэффициентами $y'=\lambda y+g(t)$ имеем $J'=0$, и член $C_2$ отсутствует. Полный список условий приведён в табл.~\ref{tab:s4cond}.

\begin{table}[ht]
\centering\small
\caption{Четырёхстадийный метод: условия и степени $h$; справа~--- показатели при $\alpha=0.3$ и $\alpha=0.7$.}\label{tab:s4cond}
\begin{tabular}{llccc}
\toprule
условие & формула & степень $h$ & $\alpha=0.3$ & $\alpha=0.7$\\
\midrule
$D_0$ & $\sum_ib_i=\gamma_0$ & $\alpha$ & 0.3 & 0.7\\
$D_1$ & $\sum_ib_ic_i=\gamma_1$ & $1+\alpha$ & 1.3 & 1.7\\
$B_1$ & $\sum_ib_i(d_1)_i=0$ & $1+2\alpha$ & 1.6 & 2.4\\
$B_2$ & $b^{\top}Ad_1=0$ & $1+3\alpha$ & 1.9 & 3.1\\
$B_3$ & $b_4a_{43}a_{32}(d_1)_2=0$ & $1+4\alpha$ & 2.2 & 3.8\\
$D_2$ & $\sum_ib_ic_i^2=\gamma_2$ & $2+\alpha$ & 2.3 & 2.7\\
$C_1$ & $\sum_ib_i(d_2)_i=0$ & $2+2\alpha$ & 2.6 & 3.4\\
$C_2$ & $\sum_ib_ic_i(d_1)_i=0$ & $2+2\alpha$ & 2.6 & 3.4\\
\bottomrule
\end{tabular}
\end{table}

\begin{theorem}[барьер для четырёх стадий]\label{thm:s4}
$q(4,\alpha)\le\max(2+2\alpha,\,1+4\alpha)$.
\end{theorem}
\begin{proof}
Все восемь показателей табл.~\ref{tab:s4cond} не превосходят $\max(2+2\alpha,1+4\alpha)$, поскольку $2+\alpha\le2+2\alpha$ и $1+3\alpha\le1+4\alpha$. Следовательно, превзойти указанное значение можно, лишь обнулив весь набор $S=\{D_0,D_1,D_2,B_1,B_2,B_3,C_1,C_2\}$. Покажем, что $S$ несовместен. По лемме~\ref{lem:chain} $B_3=b_4a_{43}a_{32}(d_1)_2$, и один из четырёх множителей должен обращаться в нуль.

(i) $b_4=0$ или $c_2=0$. Метод фактически трёхстадийный: либо последняя стадия не входит в ответ, либо вторая совпадает с первой. Набор $S$ содержит условия $D_0,D_1,D_2,B_1,B_2$ для оставшихся трёх стадий, несовместные по теореме~\ref{thm:s3}.

(ii) $a_{43}=0$ ($b_4\ne0$, $c_2>0$). Вычислим строку $b^{\top}A$: $(b^{\top}A)_3=b_4a_{43}=0$, $(b^{\top}A)_2=b_3a_{32}+b_4a_{42}$. Тогда $B_2=(b^{\top}A)_2(d_1)_2$, и поскольку $(d_1)_2\neq0$, условие $B_2=0$ даёт $b_3a_{32}+b_4a_{42}=0$. Следовательно, у строки $b^{\top}A$ отлична от нуля может быть лишь первая компонента, умножаемая на $c_1=0$, и $b^{\top}Ac=b^{\top}Ac^2=0$. Подставляя это в $B_1$ и $C_1$, находим
\[
B_1=-\frac{1}{\G(\alpha+2)}\sum_ib_ic_i^{1+\alpha},\qquad C_1=-\frac{2}{\G(\alpha+3)}\sum_ib_ic_i^{2+\alpha}.
\]
Для $\rho_i=b_ic_i$ ($i=2,3,4$) условия $D_1,D_2,B_1,C_1$ образуют систему $\sum\rho_i=\gamma_1$, $\sum\rho_ic_i=\gamma_2$, $\sum\rho_ic_i^{\alpha}=0$, $\sum\rho_ic_i^{1+\alpha}=0$ на трёх точках, т.\,е. систему (M3), несовместную по лемме~\ref{lem:measures}. Вырожденные конфигурации охватываются той же леммой: нулевой узел $c_i=0$ даёт $\rho_i=0$, а совпадающие узлы объединяются с суммарным весом, так что точек становится меньше трёх. Элементарный приём с отношением узлов здесь уже не проходит: точек три, и требуется правило знаков.

(iii) $a_{32}=0$ ($b_4\ne0$, $a_{43}\ne0$, $c_2>0$). Если $c_3=0$, то $a_{31}=0$ и $Y_3=Y_1$, т.\,е. метод фактически трёхстадийный (случай (i)); поэтому $c_3>0$. Четвёртая стадия опирается на стадии $2$ и $3$ с весами $a_{42}$, $a_{43}$; удобно рассматривать их как меру $\sigma=a_{42}\delta_{c_2}+a_{43}\delta_{c_3}$ с моментами
\[
P_1=\int x\dd\sigma=a_{42}c_2+a_{43}c_3=(Ac)_4,\qquad P_2=\int x^2\dd\sigma=a_{42}c_2^2+a_{43}c_3^2=(Ac^2)_4 .
\]
Так как $a_{32}=0$, имеем $(Ac)_2=(Ac)_3=0$ и $(Ac^2)_2=(Ac^2)_3=0$. Перепишем три условия. Поскольку $(b^{\top}A)_2=b_4a_{42}$, $(b^{\top}A)_3=b_4a_{43}$ и $(d_1)_3=-c_3^{1+\alpha}/\G(\alpha+2)$,
\[
B_2=b_4\bigl[a_{42}(d_1)_2+a_{43}(d_1)_3\bigr]=-\frac{b_4}{\G(\alpha+2)}\int x^{1+\alpha}\dd\sigma,
\]
\[
C_1=b_4P_2-\frac{2}{\G(\alpha+3)}\sum_jb_jc_j^{2+\alpha},\qquad C_2=b_4c_4P_1-\frac{1}{\G(\alpha+2)}\sum_jb_jc_j^{2+\alpha}.
\]
Итак, $B_2=0$ означает $\int x^{1+\alpha}\dd\sigma=0$, а исключение суммы $\sum_jb_jc_j^{2+\alpha}$ из $C_1=C_2=0$ с учётом $\G(\alpha+3)=(\alpha+2)\G(\alpha+2)$ даёт
\[
P_2=\frac{2c_4}{\alpha+2}\,P_1 .
\]
Если $P_1=0$, то и $P_2=0$, т.\,е. вектор $(a_{42},a_{43})$ ортогонален векторам $(c_2,c_3)$ и $(c_2^2,c_3^2)$, определитель которых $c_2c_3(c_3-c_2)$ отличен от нуля при $c_2\ne c_3$; значит, $a_{42}=a_{43}=0$, что противоречит $a_{43}\neq0$ (при $c_2=c_3$ и $a_{32}=0$ имеем $a_{31}=a_{21}$, т.\,е. $Y_3=Y_2$, и метод фактически трёхстадийный~--- случай (i)). Если $P_1\neq0$, выберем $\psi(x)=\mu_1x+\mu_2x^{1+\alpha}+\mu_3x^2\not\equiv0$ с нулями в $c_2$ и $c_3$ (три коэффициента, два условия). Тогда
\[
0=a_{42}\psi(c_2)+a_{43}\psi(c_3)=\mu_1P_1+\mu_2\cdot0+\mu_3P_2=P_1\Bigl(\mu_1+\frac{2c_4}{\alpha+2}\mu_3\Bigr),
\]
так что $\mu_1$ и $\mu_3$ имеют противоположные знаки. Функция $\psi(x)/x=\mu_1+\mu_2x^{\alpha}+\mu_3x$ с показателями $0<\alpha<1$ имеет одну перемену знака в последовательности коэффициентов и по лемме~\ref{lem:descartes}~--- не более одного положительного нуля, тогда как её нули $c_2$ и $c_3$ различны. Противоречие (при $c_4=0$ имеем $\mu_1=0$, и $\psi/x=x^{\alpha}(\mu_2+\mu_3x^{1-\alpha})$ также имеет не более одного нуля).
\end{proof}

\begin{remark}\label{rem:linear}
В случае (iii) противоречие получено из трёх условий $B_2$, $C_1$, $C_2$, причём $C_2$ в нём необходимо: численно набор $S\setminus\{C_2\}$ совместен (табл.~\ref{tab:core}). Тем самым барьер при $s=4$ порождается членом $J'g_1$, т.\,е. изменением якобиана вдоль шага, чего не было при $s=2,3$. Для линейных задач с постоянными коэффициентами ($J'=0$) этот член отсутствует, и методы, обнуляющие набор $S\setminus\{C_2\}$, имеют на таких задачах больший показатель. Построенные в п.~\ref{sec:opt} методы этим свойством не обладают (например, при $\alpha=0.3$ у них $C_1\approx-1.50$), поэтому на линейной задаче (L) их показатель равен общему значению; это и наблюдается в расчётах (табл.~\ref{tab:local}, \ref{tab:a03}).
\end{remark}

Поясним теорему на примерах (табл.~\ref{tab:s4cond}). При $\alpha=0.3$ барьер равен $\max(2.6,\,2.2)=2.6$. Чтобы его достичь, нужно обнулить всё, что лежит ниже $2.6$, в том числе член $B_3$ при $h^{2.2}$; найденный численно оптимальный метод делает это разрывом цепочки~--- у него $a_{32}=0$ (табл.~\ref{tab:coef}), а препятствием остаются $C_1,C_2$ при $h^{2.6}$. При $\alpha=0.7$ барьер равен $\max(3.4,\,3.8)=3.8$: обнуляются все члены до $C_1,C_2$ при $h^{3.4}$ включительно, препятствием служит полная цепочка $B_3$ при $h^{3.8}$, и разрыв цепочки не нужен (у оптимального метода $a_{32}\neq0$).

Численно (п.~\ref{sec:generator}) найдены четырёхстадийные методы, обнуляющие все члены с показателями меньше $\max(2+2\alpha,1+4\alpha)$, при всех двенадцати значениях $\alpha$ из табл.~\ref{tab:map} (от $0.2$ до $0.9$, включая $\frac13$ и $\frac23$); следовательно, в этих точках $q(4,\alpha)=\max(2+2\alpha,1+4\alpha)$. Сформулируем это как гипотезу.

\begin{conjecture}\label{conj4}
$q(4,\alpha)=\max(2+2\alpha,\,1+4\alpha)$ при всех $\alpha\in(0,1)$.
\end{conjecture}

\begin{remark}
При $\alpha=1$ полученные формулы дают $q=s+1$ ($s\le4$), т.\,е. классический первый барьер Бутчера $p\le s$. При $\alpha=\frac12$ они дают $q(2)=2$, $q(3)=\frac52$, $q(4)=3$, что согласуется с предложением~12 из~\cite{Lubich1983} (не существует явных двух-, трёх- и четырёхстадийных методов локального порядка $2$, $\frac52$, $3$). Формулы $1+s\alpha$ и $2+(s-2)\alpha$ совпадают ровно при $\alpha=\frac12$, поэтому излом барьера находится в этой точке. Две ветви отвечают двум стратегиям построения метода: <<цепочка>> (каждая стадия добавляет $\alpha$) и <<один разрыв цепочки>> (обнуление одной из связей $a_{i+1,i}$ даёт выигрыш на целую единицу ценой $2\alpha$).
\end{remark}

При $s\ge5$ ситуация меняется: формула $\max(1+s\alpha,\,2+(s-2)\alpha)$ при $\alpha\to1$ давала бы $q=s+1$, тогда как классический явный пятистадийный метод имеет порядок не выше четырёх, т.\,е. $q\le5$. Следовательно, вблизи $\alpha=1$ должен проявиться дробный аналог так называемых <<глубоких>> барьеров Бутчера. Численные результаты п.~\ref{sec:s5} приводят к следующей гипотезе.

\begin{conjecture}\label{conj5}
Для методов класса~\eqref{eq:rowsum} с $\hat c=1.5$
\[
q(5,\alpha)=\max\Bigl(2+3\alpha,\;\min\bigl(1+5\alpha,\;3+2\alpha\bigr)\Bigr),\qquad 0<\alpha<1 .
\]
\end{conjecture}

Барьер для пяти стадий имеет, таким образом, два излома: при $\alpha=\frac12$, как и раньше, и при $\alpha=\frac23$, где <<цепочечная>> ветвь $1+5\alpha$ срезается ветвью $3+2\alpha$. При $\alpha=1$ последняя даёт ровно классическое значение $q=5$, так что дробный барьер непрерывно сшивается с классическим.

Удобная формулировка для общего случая такова. Каждой стадии сопоставим меру $\sigma_i=\sum_ja_{ij}\delta_{c_j}$ и оператор выталкивания $T\mu=\sum_i\mu_i\sigma_i$ (явность означает $T^s=0$); точному решению отвечает оператор $T^{*}$ с ядром $(c-\theta)^{\alpha-1}/\G(\alpha)$. Условия, препятствующие превышению значения $\max(1+s\alpha,\,2+(s-2)\alpha)$, означают совпадение $T$ и $T^{*}$ на моментах $\theta$ и $\theta^2$ вдоль орбиты $b,Tb,T^2b,\dots$ и на моменте $\theta$ вдоль орбит мер $c\cdot T^jb$. Индукция по $s$ закрывает случаи $b_s=0$ и $c_2=0$; случай разрыва $a_{i+1,i}=0$ при общем $s$, как и строгое описание третьей ветви, остаются открытыми.
\subsection{Вычисление лаг-члена}\label{sec:lag}

На практике $\Phi_n$ вычисляется по уже найденным значениям. В классических методах Пуазе используются значения стадий $F_{m,j}$ на прошлых шагах.

\begin{claim}[численное наблюдение]\label{cl:stagelag}
Если $\Phi_n$ вычисляется произведенческой интерполяцией значений стадий $F_{m,j}$ на прошлых интервалах, то при $\alpha<1$ глобальная ошибка явного метода имеет, вообще говоря, порядок $h^{1+\alpha}$ независимо от числа стадий.
\end{claim}
\noindent\emph{Обоснование.} Ошибки стадий равны $e_{m,j}=h^{1+\alpha}g_1(d_1)_j+\dots$, причём $(d_1)_2=-c_2^{1+\alpha}/\G(\alpha+2)\neq0$. Они входят в $\Phi_n$ с весами $h^{\alpha}W_j\bigl((t-t_m)/h\bigr)$, сумма которых по $m$ ограничена, что даёт вклад $O(h^{1+\alpha})$. Взаимное уничтожение этих вкладов потребовало бы, чтобы многочлен, интерполирующий дефекты $(d_1)_j$ по узлам $c_j$, был ортогонален ядру $(x-\theta)^{\alpha-1}$ при всех $x\ge1$, что при $\alpha<1$ невозможно для ненулевого многочлена. При $\alpha=1$ ядро постоянно и такое уничтожение возможно: одна и та же квадратура $b$ годится для любой целевой точки, и именно поэтому классический метод Рунге--Кутты может оставаться одношаговым. Это рассуждение не является доказательством: строгий анализ требует оценки остатков и здесь не проводится, так что утверждение~\ref{cl:stagelag} следует рассматривать как численно подтверждённое наблюдение (рис.~\ref{fig:lag}), а не как теорему.

Поэтому лаг-член вычисляется по принятым значениям $F_m=f(t_m,y_m)$ в узлах сетки, интерполируя их на каждом прошлом интервале многочленом степени $r$ по $r+1$ соседним узлам и интегрируя с ядром точно. Погрешность такой квадратуры для $g\in C^{r+1}[0,T]$ имеет порядок $O(h^{r+1})$~\cite{Brunner2004,Diethelm2010}. Решения уравнений с производной Капуто, однако, вообще говоря негладки в нуле: $g$ содержит члены вида $t^{j+k\alpha}$~\cite{Lubich1983}. Для тестовых задач п.~\ref{sec:global} младший негладкий член $g$ имеет вид $t^{4-\alpha}$ (соответственно $t^{5-\alpha}$), т.\,е. $g\in C^4$, но $g\notin C^5$. Погрешность интерполяции такого члена на интервале $[t_k,t_{k+1}]$ равна $O\bigl(h^{r+1}(kh)^{3-\alpha-r}\bigr)$ при $k\ge1$ и $O(h^{4-\alpha})$ на первых интервалах; при ограниченном ядре суммирование даёт вклад $O(h^{5-\alpha})$, а вблизи текущей точки $g$ гладкая. Тем самым погрешность квадратуры для этих задач имеет порядок $O(h^{\min(r+1,\,5-\alpha)})$, что при $r=4$ выше показателя $q\le3.8$ всех рассматриваемых методов. Если же $g\sim t^{\alpha}$, как в общем случае, погрешность первых интервалов становится главной, и для сохранения порядка нужны сгущающиеся сетки или стартовые поправки~\cite{Lubich1983}; этот случай здесь не рассматривается. Утверждение о том, что при достаточно точной квадратуре глобальная ошибка определяется показателем $q$, в работе не доказывается: оно подтверждено численно, а строгий анализ накопления ошибок требует отдельного исследования. Отметим, что глобальный порядок $q$, а не $q-1$, как для обыкновенных уравнений,~--- свойство не дробного случая, а именно такого способа вычисления лаг-члена: ошибка, допущенная на шаге, не переносится целиком на следующие шаги, а входит в них лишь через значение $f(t_{n+1},y_{n+1})$ с весом порядка $h^{\alpha}$. Тот же механизм при $\alpha=1$ дал бы и для классического метода Рунге--Кутты глобальный порядок $p+1$. Иными словами, такая реализация~--- это гибрид многошаговой квадратуры предыстории и рунге--куттовского предиктора нового значения, и барьеры раздела~2 ограничивают именно предиктор.

\section{Алгоритм и результаты}

\subsection{Генератор условий порядка}\label{sec:generator}

Условия порядка строятся автоматически. Локальная ошибка при точном лаг-члене разлагается в ряд по двум переменным $H_1=h^{\alpha}$ и $H=h$; коэффициенты являются многочленами от формальных атомов $g_k=\frac{\dd^k}{\dd t^k}f(t,y(t))$ и $D_{r,k}=\frac{\dd^k}{\dd t^k}\partial_y^rf(t,y(t))$ в точке $t_n$. Алгоритм состоит из следующих шагов:
\begin{enumerate}
\item рекурсивно по стадиям вычислить ряды для ошибок стадий $e_i$ и значений $f(t_n+c_ih,y(t_n+c_ih)+e_i)$, отбрасывая мономы $H_1^mH^k$ с $m\alpha+k>p$;
\item сгруппировать мономы с одинаковым значением $m\alpha+k$ (это существенно при рациональных $\alpha$ с малым знаменателем);
\item каждое ненулевое произведение атомов при показателе $\le p$ даёт одно условие порядка;
\item существование метода с показателем больше $p$ проверить минимизацией суммы квадратов невязок (метод Левенберга--Марквардта из 7--60 случайных начальных точек, $0\le c_i\le\hat c$): невязка $\sim10^{-15}$ означает существование, устойчиво отделённая от нуля невязка ($10^{-6}$--$10^{-2}$) свидетельствует о несуществовании.
\end{enumerate}
\paragraph{Реализация.} Генератор написан на языке Python с использованием библиотек NumPy, SciPy и SymPy; все расчёты выполнялись в двойной точности. Ряды хранятся как словари, сопоставляющие паре показателей $(m,k)$ и мономиалу атомов числовой (в символьной версии~--- символьный) коэффициент. Значение $\alpha$ приводится к рациональному виду с знаменателем не больше $1000$, что обеспечивает точную группировку мономов с равными показателями. Поиск методов выполняется функцией \texttt{scipy.optimize.least\_squares} (метод доверительной области с ограничениями $0\le c_i\le\hat c$) с допусками \texttt{xtol}${}=$\texttt{ftol}${}=$\texttt{gtol}${}=10^{-15}$ и не более 3000--4000 вычислениями невязки на старт; методом считается найденным, если невязка не превосходит $10^{-10}$. Подчеркнём, что неудача поиска (невязка $10^{-6}$--$10^{-2}$) является лишь численным свидетельством несуществования решения, но не его доказательством: она может быть следствием плохой обусловленности или узкой области существования решений. Именно поэтому все утверждения о несуществовании при $s\ge5$ сформулированы в работе как гипотезы. Код генератора, расчётные скрипты и данные, по которым построены все таблицы и рисунки, открыты в репозитории \url{https://github.com/PrikazchikovDA/fractional-vrk-barriers}.

Генератор проверен на классическом случае $\alpha=1$ (максимальный порядок $s$ при $s=2,3,4$) и на результате Любича при $\alpha=\frac12$ (предложение~12 и пример~13 из~\cite{Lubich1983}). Для анализа барьера при $s\ge5$ использовалось также ядро условий, записанное в матричной форме (условия $D_0,D_1,D_2$, цепочка $B_m$, цепочка $b^{\top}A^qd_2$ и вставки $b^{\top}A^j\mathrm{diag}(c)A^{q-j}d_1$), что сокращает вычисления на порядки.

\subsection{Карта барьеров}\label{sec:map}

Результаты расчёта показателей приведены в табл.~\ref{tab:map} и на рис.~\ref{fig:barriers}. При $s\le4$ во всех точках численное значение совпадает с формулой $\max(1+s\alpha,\,2+(s-2)\alpha)$, при $s=5$~--- с формулой гипотезы~\ref{conj5}.

\begin{table}[ht]
\centering\small
\caption{Наилучший показатель локальной ошибки $q(s,\alpha)$: численный расчёт. Для $s=5$ узлы ограничены условием $c_i\le1.5$.}\label{tab:map}
\begin{tabular}{lcccccccccccc}
\toprule
$\alpha$ & 0.2 & 0.25 & 0.3 & $\frac13$ & 0.4 & 0.5 & 0.6 & $\frac23$ & 0.7 & 0.75 & 0.8 & 0.9\\
\midrule
$s=2$ & 1.4 & 1.5 & 1.6 & 1.67 & 1.8 & 2.0 & 2.2 & 2.33 & 2.4 & 2.5 & 2.6 & 2.8\\
$s=3$ & 2.2 & 2.25 & 2.3 & 2.33 & 2.4 & 2.5 & 2.8 & 3.0 & 3.1 & 3.25 & 3.4 & 3.7\\
$s=4$ & 2.4 & 2.5 & 2.6 & 2.67 & 2.8 & 3.0 & 3.4 & 3.67 & 3.8 & 4.0 & 4.2 & 4.6\\
$s=5$ & 2.6 & 2.75 & 2.9 & 3.0 & 3.2 & 3.5 & 4.0 & 4.33 & 4.4 & 4.5 & 4.6 & 4.8\\
\midrule
формула & \multicolumn{12}{l}{$s=2$: $1+2\alpha$;\quad $s=3,4$: $\max(1+s\alpha,\,2+(s-2)\alpha)$;\quad $s=5$: гипотеза~\ref{conj5}}\\
\bottomrule
\end{tabular}
\end{table}

\begin{figure}[ht]
\centering
\includegraphics[width=0.66\textwidth]{figs/barriers.pdf}
\caption{Дробный первый барьер: линии~--- формулы (для $s=5$~--- гипотеза~\ref{conj5}), маркеры~--- численный расчёт, звёздочки~--- классические значения при $\alpha=1$ ($q=s+1$ при $s\le4$ и $q=5$ при $s=5$). Тонкая точечная линия~--- формальное продолжение ветви $1+5\alpha$.}\label{fig:barriers}
\end{figure}

\begin{table}[ht]
\centering\small
\caption{Минимальная невязка ядра условий при попытке превысить значение $\max(1+s\alpha,\,2+(s-2)\alpha)$. <<Векторный>> вариант~--- условия для систем, <<скалярный>>~--- вставки $J'$ просуммированы. Невязка $\lesssim10^{-13}$ означает совместность.}\label{tab:core}
\begin{tabular}{ccccccc}
\toprule
$s$ & $\alpha$ & векторный & скалярный & без вставок $J'$ & без цепочки $g_2$ & вывод\\
\midrule
3 & 0.3 & $9\cdot10^{-4}$ & $9\cdot10^{-4}$ & $9\cdot10^{-4}$ & $9\cdot10^{-4}$ & несовместно\\
3 & 0.6 & $3\cdot10^{-5}$ & $3\cdot10^{-5}$ & $3\cdot10^{-5}$ & $3\cdot10^{-5}$ & несовместно\\
4 & 0.3 & $5\cdot10^{-3}$ & $5\cdot10^{-3}$ & $5\cdot10^{-14}$ & $2\cdot10^{-16}$ & несовместно\\
4 & 0.6 & $4\cdot10^{-5}$ & $4\cdot10^{-5}$ & $7\cdot10^{-16}$ & $5\cdot10^{-13}$ & несовместно\\
5 & 0.3 & $1\cdot10^{-3}$ & $7\cdot10^{-4}$ & $3\cdot10^{-14}$ & $3\cdot10^{-16}$ & несовместно\\
5 & 0.6 & $3\cdot10^{-6}$ & $2\cdot10^{-5}$ & $3\cdot10^{-16}$ & $5\cdot10^{-15}$ & несовместно\\
6 & 0.3 & $7\cdot10^{-5}$ & $3\cdot10^{-6}$ & $2\cdot10^{-15}$ & $7\cdot10^{-16}$ & вероятно несовм.\\
6 & 0.6 & $1\cdot10^{-6}$ & $9\cdot10^{-16}$ & $1\cdot10^{-15}$ & $6\cdot10^{-16}$ & скаляр: совместно\\
\bottomrule
\end{tabular}
\end{table}

Табл.~\ref{tab:core} показывает два эффекта. Во-первых, при $s\ge4$ барьер держится одновременно на цепочке членов с $g_2$ и на членах с $J'$: удаление любой из групп делает ядро совместным. Во-вторых, при $s=6$ скалярный и векторный варианты расходятся (для скалярных уравнений ядро при $\alpha=0.6$ совместно), аналогично тому, как в классической теории скалярные и векторные условия порядка различаются начиная с порядка~5~\cite{Butcher2016}. Поэтому все гипотезы сформулированы для систем.

\subsection{Пять стадий: третья ветвь барьера}\label{sec:s5}

Рассмотрим подробнее случай $s=5$, где формула $\max(1+5\alpha,\,2+3\alpha)$ вступает в противоречие с классикой. Результаты поиска методов с полным набором условий порядка (скалярный генератор, от 27 до 90 условий на 15 неизвестных в зависимости от $\alpha$ и уровня) приведены в табл.~\ref{tab:s5}.

\begin{table}[ht]
\centering\small
\caption{Поиск пятистадийных методов, обнуляющих все члены с показателями до указанного (включительно), при $c_i\le1.5$. Для строк <<нет>> указана наименьшая невязка по всем случайным стартам. Последний столбец~--- максимальный модуль коэффициента найденного метода.}\label{tab:s5}
\begin{tabular}{cccccc}
\toprule
$\alpha$ & условия до & найден & невязка & $\max|\text{коэф.}|$ & вывод\\
\midrule
0.2 & 2.4 & да & $9\cdot10^{-16}$ & 2.7 & \multirow{2}{*}{$q=2.6=2+3\alpha$}\\
0.2 & 2.6 & нет & $1\cdot10^{-2}$ & --- & \\
0.25 & 2.5 & да & $4\cdot10^{-16}$ & 3.6 & \multirow{2}{*}{$q=2.75=2+3\alpha$}\\
0.25 & 2.75 & нет & $2\cdot10^{-2}$ & --- & \\
$\frac13$ & $\frac83$ & да & $3\cdot10^{-16}$ & 2.7 & \multirow{2}{*}{$q=3=2+3\alpha$}\\
$\frac13$ & 3 & нет & $5\cdot10^{-3}$ & --- & \\
0.4 & 3.0 & да & $3\cdot10^{-16}$ & 3.7 & \multirow{2}{*}{$q=3.2=2+3\alpha$}\\
0.4 & 3.2 & нет & $6\cdot10^{-4}$ & --- & \\
0.5 & 3.0 & да & $2\cdot10^{-15}$ & 4.0 & \multirow{2}{*}{$q=3.5=2+3\alpha=1+5\alpha$}\\
0.5 & 3.5 & нет & $1\cdot10^{-4}$ & --- & \\
0.62 & 3.86 & да & $1\cdot10^{-15}$ & 6.4 & $q\ge4.1=1+5\alpha$\\
$\frac23$ & 4.0 & да & $6\cdot10^{-16}$ & 5.3 & \multirow{2}{*}{$q=\frac{13}{3}=1+5\alpha=3+2\alpha$}\\
$\frac23$ & $\frac{13}{3}$ & нет & $1\cdot10^{-5}$ & --- & \\
0.7 & 4.2 & да & $7\cdot10^{-16}$ & 4.2 & \multirow{2}{*}{$q=4.4=3+2\alpha$}\\
0.7 & 4.4 & нет & $4\cdot10^{-6}$ & --- & \\
0.75 & 4.25 & да & $4\cdot10^{-16}$ & 4.9 & \multirow{2}{*}{$q=4.5=3+2\alpha$}\\
0.75 & 4.5 & нет & $3\cdot10^{-6}$ & --- & \\
0.8 & 4.4 & да & $3\cdot10^{-16}$ & 2.4 & \multirow{2}{*}{$q=4.6=3+2\alpha$}\\
0.8 & 4.6 & нет & $3\cdot10^{-6}$ & --- & \\
0.9 & 4.7 & да & $2\cdot10^{-16}$ & 2.9 & \multirow{2}{*}{$q=4.8=3+2\alpha$}\\
0.9 & 4.8 & нет & $4\cdot10^{-6}$ & --- & \\
\bottomrule
\end{tabular}
\end{table}

Как видно, при $\alpha\le\frac12$ показатель определяется ветвью $2+3\alpha$, при $\alpha=\frac23$ ветви $1+5\alpha$ и $3+2\alpha$ встречаются, а при $\alpha\ge0.7$ показатель определяется ветвью $3+2\alpha$. Во всех рассмотренных значениях $\alpha$ результаты расчёта согласуются с формулой гипотезы~\ref{conj5}, причём значение $q(5;0.9)=4.8$ было предсказано формулой гипотезы~\ref{conj5} до проведения расчёта. Коэффициенты найденных методов имеют порядок единиц, т.\,е. речь идёт о вполне <<нормальных>> методах.

Интересно, что без ограничения на узлы картина несколько иная. При $\alpha=0.7$ находится метод, обнуляющий все члены до показателя $4.4$, однако он имеет узлы $c\approx(0;\,0.48;\,9.39;\,0.73;\,4.49)$, веса $b\approx(0.21;\,0;\,0;\,0.89;\,0.001)$ и почти взаимно уничтожающиеся коэффициенты $a_{31}\approx-3188$, $a_{32}\approx3193$. Продолжение этого решения по параметру $\alpha$ (табл.~\ref{tab:path}) показывает, что оно образует гладкую ветвь, коэффициенты на которой неограниченно возрастают при $\alpha\to\alpha^*\approx0.72$, после чего ветвь исчезает. Тем самым формально ветвь $1+5\alpha$ продолжается за точку $\alpha=\frac23$, но лишь за счёт патологических методов, вычисляющих правую часть на 5--10 шагов вперёд и обладающих коэффициентами порядка $10^3$--$10^4$. Практического интереса такие методы не представляют, что и оправдывает ограничение $c_i\le\hat c$ в определении класса~\eqref{eq:rowsum}.

\begin{table}[ht]
\centering\small
\caption{Продолжение по $\alpha$ патологического пятистадийного метода (все члены с показателями меньше $1+5\alpha$ обнулены, невязка $\le3\cdot10^{-14}$; при $\alpha=0.716$ итерации не сходятся).}\label{tab:path}
\begin{tabular}{lcccccccc}
\toprule
$\alpha$ & 0.68 & 0.69 & 0.70 & 0.704 & 0.708 & 0.712 & 0.714 & 0.716\\
\midrule
$\max|\text{коэф.}|$ & 1300 & 1890 & 3190 & 4250 & 6160 & 10750 & 17370 & ---\\
\bottomrule
\end{tabular}
\end{table}

\subsection{Оптимальные методы}\label{sec:opt}

Трёхстадийные методы задаются семействами A и B из доказательства теоремы~\ref{thm:s3}; в расчётах использованы узлы $c=(0,\frac12,1)$. Найденные четырёхстадийные методы приведены в табл.~\ref{tab:coef}. Точность записи коэффициентов существенна: при округлении до шести знаков невязки условий возрастают до $10^{-7}$--$10^{-6}$, что на мелких сетках может исказить наблюдаемую асимптотику (невязка всех условий при исходных коэффициентах $\sim10^{-16}$, узлы подбирались в $(0,1.5]$ с минимизацией максимального коэффициента).

\begin{table}[ht]
\centering
\caption{Четырёхстадийные методы с показателем $q(4,\alpha)$ в виде таблиц Бутчера: в левом столбце узлы $c_i$, справа от черты матрица $A$, под чертой веса $b_j$. Значения округлены до шести знаков после запятой; с полной точностью коэффициенты приведены в файле \texttt{data/tab4.json} репозитория, и для воспроизведения порядков на мелких сетках следует использовать именно их.}\label{tab:coef}
{\footnotesize\setlength{\arraycolsep}{3pt}%
\begin{minipage}[t]{0.49\textwidth}\centering
$\alpha=0.3$,\quad $q=2+2\alpha=2.6$\\[4pt]
$\begin{array}{r|rrrr}
0 &  &  &  & \\
0.390638 & 0.840451 &  &  & \\
0.886375 & 1.074645 & 0 &  & \\
0.996293 & -0.717071 & 2.792600 & -0.962527 & \\
\hline
 & -0.439290 & 1.677188 & -2.957963 & 2.834308
\end{array}$
\end{minipage}\hfill
\begin{minipage}[t]{0.49\textwidth}\centering
$\alpha=0.7$,\quad $q=1+4\alpha=3.8$\\[4pt]
$\begin{array}{r|rrrr}
0 &  &  &  & \\
0.684259 & 0.843846 &  &  & \\
0.847616 & 1.874949 & -0.894673 &  & \\
1.070235 & 1.385151 & 0.530238 & -0.761288 & \\
\hline
 & 0.169650 & 1.307226 & -0.699223 & 0.322894
\end{array}$
\end{minipage}}
\end{table}

\subsection{Глобальная сходимость}\label{sec:global}

Перейдём к проверке полученных методов на задачах с известными точными решениями ($T=1$):
\begin{itemize}
\item[(L)] линейная задача из~\cite{Ghaffari2023}: $D^{\alpha}y=-y+t^{4-\alpha}/\G(5-\alpha)$, $y=t^4E_{\alpha,5}(-t^{\alpha})$;
\item[(N)] нелинейное уравнение $D^{\alpha}y=\frac{24}{\G(5-\alpha)}t^{4-\alpha}+y^2-t^8$, $y=t^4$;
\item[(S)] неавтономная нелинейная система $D^{\alpha}y_1=\frac{24\,t^{4-\alpha}}{\G(5-\alpha)}+y_1y_2-t^9-\sin y_2+\sin t^5$, $D^{\alpha}y_2=\frac{120\,t^{5-\alpha}}{\G(6-\alpha)}-y_1^2+t^8+t(y_1-t^4)$, $y=(t^4,t^5)$.
\end{itemize}
Лаг-член вычислялся по значениям в узлах сетки с интерполяцией степени $r=4$; стартовые значения $y_0,\dots,y_4$ брались точными, чтобы отделить порядок шага от стартовой процедуры. Порядок оценивался как $\log_2\bigl(E(N)/E(2N)\bigr)$ для ошибки $E(N)=\|y_N-y(1)\|_\infty$.

Прежде чем обсуждать глобальную сходимость, проверим показатели построенных методов напрямую: сделаем один шаг из точки $t_n=0.5$ с точным лаг-членом (интеграл $\Phi_n$ вычислялся адаптивной квадратурой с весом $(t-s)^{\alpha-1}$ с относительной точностью $10^{-14}$) и оценим показатель как $\log_2\bigl(\delta(h)/\delta(h/2)\bigr)$. Результаты приведены в табл.~\ref{tab:local}.

\begin{table}[ht]
\centering\small
\caption{Показатель локальной ошибки, измеренный напрямую (шаг из $t_n=0.5$ с точным лаг-членом, $h=1/80\to1/160\to1/320$).}\label{tab:local}
\begin{tabular}{lcccc}
\toprule
метод & прогноз $q$ & (L) & (N) & (S)\\
\midrule
\multicolumn{5}{l}{$\alpha=0.7$}\\
2 стадии & $1+2\alpha=2.4$ & 2.46, 2.45 & 3.07, 3.40 & 2.49, 2.47\\
3 стадии, A & $2+\alpha=2.7$ & 2.72, 2.71 & 2.71, 2.70 & 2.72, 2.71\\
3 стадии, B & $1+3\alpha=3.1$ & 3.06, 3.07 & 3.56, 3.53 & 3.25, 3.21\\
4 стадии & $1+4\alpha=3.8$ & 3.89, 3.87 & 4.14, 4.12 & 4.12, 4.10\\
\midrule
\multicolumn{5}{l}{$\alpha=0.3$}\\
2 стадии & $1+2\alpha=1.6$ & 1.63, 1.62 & 1.59, 1.59 & 1.64, 1.62\\
3 стадии, A & $2+\alpha=2.3$ & 2.35, 2.34 & 2.29, 2.29 & 2.34, 2.33\\
3 стадии, B & $1+3\alpha=1.9$ & 1.88, 1.89 & 1.90, 1.89 & 1.98, 1.95\\
4 стадии & $2+2\alpha=2.6$ & 2.65, 2.64 & 2.67, 2.65 & 2.59, 2.58\\
\bottomrule
\end{tabular}
\end{table}

При $\alpha=0.3$ все измеренные значения совпадают с прогнозом с точностью до нескольких сотых, в том числе для четырёхстадийного метода на линейной задаче: у этого метода $C_1\approx-1.50\ne0$, и препятствующий член $h^{2+2\alpha}Jg_2$ присутствует и при $J'=0$ (замечание~\ref{rem:linear}). При $\alpha=0.7$ на линейной задаче прогноз также подтверждается, а на нелинейных задачах часть показателей выше прогноза. Это не противоречит теории: $q$~--- оценка для \emph{всех} гладких $f$, тогда как для конкретной задачи коэффициент при главном члене может оказаться малым. Препятствующий член четырёхстадийного метода при $\alpha=0.7$~--- полная цепочка $h^{1+4\alpha}B_3J^3g_1$ с $B_3\approx-0.075$, т.\,е. коэффициент пропорционален $\|J\|^3$. Для задачи (N) имеем $J=2y=2t^4$, и при $t\le1$ величина $|J|^3$ мала (в точке $t=0.5$ она равна $0.125^3$); для системы (S) $\|J^3\|\approx0.33$. Поэтому на доступных шагах главным оказывается следующий член разложения, и наблюдаемый показатель близок к $4.1$. Аналогично для двухстадийного метода и метода B на задаче (N): главный член содержит множитель $J$ (соответственно $J^2$), а измерение в других точках показывает, что коэффициент при нём меняется вдоль решения и проходит через нуль (при $t_n=0.25$ показатель двухстадийного метода убывает $2.75\to2.71$, при $t_n=0.75$ он равен $2.36$, а для метода B при $t_n=0.75$ наблюдается характерное для сокращения членов немонотонное поведение). Таким образом, превышение прогноза здесь является свойством задачи, а не метода.

Перейдём к глобальной сходимости (табл.~\ref{tab:a07}, \ref{tab:a03}).

\begin{table}[ht]
\centering\small
\caption{Наблюдаемые глобальные порядки при $\alpha=0.7$, $N=80\to160\to320\to640\to1280$.}\label{tab:a07}
\begin{tabular}{lcccc}
\toprule
метод & прогноз $q$ & (L) & (N) & (S)\\
\midrule
2 стадии & $1+2\alpha=2.4$ & 2.44, 2.43, 2.42, 2.42 & 2.31, 2.34, 2.36, 2.38 & 2.35, 2.37, 2.38, 2.38\\
3 стадии, A & $2+\alpha=2.7$ & 2.73, 2.72, 2.71, 2.71 & 2.55, 2.61, 2.65, 2.67 & 2.74, 2.73, 2.72, 2.71\\
3 стадии, B & $1+3\alpha=3.1$ & 3.10, 3.10, 3.10, 3.10 & 2.90, 2.97, 3.01, 3.04 & 3.81, 2.86, 2.87, 2.94\\
4 стадии & $1+4\alpha=3.8$ & 3.85, 3.84, 3.83, 3.83 & 3.89, 3.90, 3.90, 3.89 & 4.06, 4.08, 4.09, 4.09\\
\bottomrule
\end{tabular}
\end{table}

\begin{table}[ht]
\centering\small
\caption{Наблюдаемые глобальные порядки при $\alpha=0.3$, $N=160\to320\to640\to1280\to2560$.}\label{tab:a03}
\begin{tabular}{lcccc}
\toprule
метод & прогноз $q$ & (L) & (N) & (S)\\
\midrule
2 стадии & $1+2\alpha=1.6$ & 1.65, 1.64, 1.64, 1.63 & 1.44, 1.49, 1.51, 1.53 & 1.54, 1.55, 1.57, 1.57\\
3 стадии, A & $2+\alpha=2.3$ & 2.30, 2.32, 2.32, 2.33 & 1.56, 1.79, 1.93, 2.02 & 2.38, 2.41, 2.42, 2.42\\
3 стадии, B & $1+3\alpha=1.9$ & 1.94, 1.94, 1.94, 1.93 & 1.72, 1.77, 1.80, 1.82 & 1.80, 1.84, 1.86, 1.87\\
4 стадии & $2+2\alpha=2.6$ & 2.70, 2.68, 2.66, 2.65 & 2.04, 2.19, 2.29, 2.36 & 2.36, 2.41, 2.45, 2.48\\
\bottomrule
\end{tabular}
\end{table}

\begin{figure}[ht]
\centering
\includegraphics[width=\textwidth]{figs/convergence_a07.pdf}
\caption{Сходимость при $\alpha=0.7$; тонкие линии~--- наклоны $N^{-q}$ с теоретическими показателями.}\label{fig:conv}
\end{figure}

При $\alpha=0.7$ (табл.~\ref{tab:a07}, рис.~\ref{fig:conv}) для линейной задачи наблюдаемые порядки совпадают с прогнозом с точностью до сотых; для задачи (N) они монотонно приближаются к прогнозу снизу. Превышение прогноза четырёхстадийным методом на (N) и (S) объясняется малостью коэффициента $B_3J^3g_1$, о чём говорилось выше. Для метода B на системе (S) первый порядок $3.81$ выпадает из ряда: на грубой сетке два старших члена ошибки частично взаимно уничтожаются, и на последовательном измельчении ($N$ до $1280$, ошибка до $4\cdot10^{-10}$) порядки $2.86, 2.87, 2.94$ монотонно возрастают к прогнозу $3.1$. При $\alpha=0.3$ (табл.~\ref{tab:a03}) выход на асимптотику медленнее, поскольку соседние члены разложения различаются лишь множителем $h^{0.3}$; тем не менее на данном диапазоне сеток все последовательности порядков монотонны и приближаются к прогнозу. Для четырёхстадийного метода на линейной задаче порядки убывают $2.70\to2.65$ к значению $2.6$, что согласуется с прямым измерением локального показателя ($2.64$). Метод A на системе (S) даёт $2.42$ при прогнозе $2.3$ и локальном показателе $2.33$: здесь, как и выше, младший член глобальной ошибки для данной задачи мал, и наблюдаемый порядок определяется на доступных сетках следующим членом. Подчеркнём, что глобальные расчёты выполнены при двух значениях $\alpha=0.3$ и $0.7$; они подтверждают теоретические показатели конкретных методов на выбранных задачах, но не являются эмпирическим обобщением на весь интервал $(0,1)$~--- для этого служат теоремы~\ref{thm:s2}--\ref{thm:s4} и карта барьеров п.~\ref{sec:map}.

\begin{figure}[ht]
\centering
\includegraphics[width=0.78\textwidth]{figs/lag.pdf}
\caption{Влияние способа вычисления лаг-члена (задача (L)): при интерполяции по значениям стадий порядок ограничен величиной $\approx1+\alpha$ (утверждение~\ref{cl:stagelag}), при интерполяции по узлам сетки достигается показатель $q$.}\label{fig:lag}
\end{figure}

\subsection{Проверка методов EFORK}

Реализация двухстадийного метода EFORK воспроизводит таблицы 1--2 работы~\cite{Ghaffari2023} во всех приведённых знаках (например, $E=2.05503\cdot10^{-3}$ при $\alpha=\frac12$, $h=\frac1{40}$ и $E=1.09027\cdot10^{-2}$ при $\alpha=\frac13$, $h=\frac1{40}$). Результаты при других $\alpha$ приведены в табл.~\ref{tab:efork} и на рис.~\ref{fig:efork}.

\begin{table}[ht]
\centering\small
\caption{Методы EFORK~\cite{Ghaffari2023} на задаче (L). $R$~--- коэффициент из теоремы~\ref{thm:efork}; <<лок.>>~--- порядок локальной ошибки при точной истории (шаг из $t=0.5$, $h\to0$); <<глоб.>>~--- наблюдаемый глобальный порядок ($N=160\to640$). Трёхстадийный метод взят с $c=(0,\frac12,1)$.}\label{tab:efork}
\begin{tabular}{cccccc}
\toprule
$\alpha$ & $R$ (2 ст.) & лок. (2 ст.) & глоб. (2 ст.), заявл. $2\alpha$ & $R$ (3 ст.) & глоб. (3 ст.), заявл. $3\alpha$\\
\midrule
0.25 & 3.826 & 1.18 & 1.00 \;(0.50) & 0.996 & 1.77 \;(0.75)\\
$\frac13$ & 1.936 & 1.16 & 0.88 \;(0.67) & 1.000 & 1.94 \;(1.00)\\
0.4 & 1.386 & 1.19 & 0.83 \;(0.80) & 1.002 & 2.07 \;(1.20)\\
0.5 & 1.000 & 1.61 & 1.10 \;(1.00) & 1.000 & 1.56 \;(1.50)\\
0.6 & 0.808 & 0.83 & 0.20 \;(1.20) & 0.990 & 0.37 \;(1.80)\\
0.7 & 0.693 & 0.95 & 0.23 \;(1.40) & 0.969 & 0.26 \;(2.10)\\
0.8 & 0.614 & 0.98 & 0.15 \;(1.60) & 0.936 & 0.17 \;(2.40)\\
0.9 & 0.553 & 0.99 & 0.07 \;(1.80) & 0.891 & 0.08 \;(2.70)\\
\bottomrule
\end{tabular}
\end{table}

\begin{figure}[ht]
\centering
\includegraphics[width=\textwidth]{figs/efork.pdf}
\caption{(а) Коэффициент согласованности $R(\alpha)$ методов EFORK; согласованность требует $R=1$. (б) Наблюдаемые глобальные порядки в сравнении с заявленными $2\alpha$ и $3\alpha$.}\label{fig:efork}
\end{figure}

По представленным данным отчётливо видно, что при $R\ne1$ локальная ошибка имеет порядок $h^1$ (столбец <<лок.>>), и заявленный порядок не достигается. Значения в столбцах <<глоб.>> получены на диапазоне $N=160\to640$ и характеризуют предасимптотический режим, а не асимптотический порядок: при $\alpha\ge0.6$ наблюдаемый порядок близок к нулю, а при малых $\alpha$ он близок к единице~--- иногда даже выше заявленного $2\alpha$, но определяется уже не условиями порядка метода, а его несогласованностью. При $\alpha\le0.4$ трёхстадийный метод имеет $|R-1|\le4\cdot10^{-3}$, поэтому несогласованность проявляется лишь на мелких сетках: при $\alpha=0.4$ ошибка перестаёт убывать на уровне $\approx5\cdot10^{-7}$ при $N=2560$.

\section{Выводы}

Результатом настоящей работы является описание барьеров порядка явных одношаговых методов типа Рунге--Кутты для дробно-дифференциальных уравнений с производной Капуто. Основные результаты состоят в следующем.
\begin{enumerate}
\item \emph{Доказано}, что необходимым условием согласованности методов Рунге--Кутты с переносом базовой точки производной Капуто (EFORK) на внутренних шагах является равенство $\sum_iw_ic_i^{1-\alpha}=\G(2-\alpha)$, которое следует из их условий порядка лишь для конечного набора значений $\alpha=1/m$; именно при таких $\alpha$ проводились численные эксперименты исходной работы. Для опубликованных методов это условие при прочих $\alpha$ нарушено, что подтверждено и прямым расчётом.
\item Для явных методов в форме Вольтерры \emph{доказано} представление младших членов локальной ошибки через квадратурные дефекты $D_k$, цепочки $B_m=b^{\top}A^{m-1}d_1$ и члены $C_1,C_2$, связанные с $f_yg_2$ и изменением якобиана вдоль решения.
\item \emph{Доказан} дробный аналог первого барьера Бутчера: $q(2,\alpha)=1+2\alpha$, $q(3,\alpha)=\max(2+\alpha,1+3\alpha)$, $q(4,\alpha)\le\max(2+2\alpha,1+4\alpha)$. Достижимость оценки при $s=4$ установлена \emph{численно} в двенадцати значениях $\alpha$ (построены соответствующие методы); равенство при всех $\alpha$ остаётся \emph{гипотезой}~\ref{conj4}. Ключевой инструмент~--- правило знаков Декарта для мер, сосредоточенных в узлах. Барьер имеет излом при $\alpha=\frac12$ и при $\alpha=1$ переходит в классическое ограничение $p\le s$.
\item При пяти стадиях расчёты приводят к \emph{гипотезе}~\ref{conj5} о третьей ветви барьера: $q(5,\alpha)=\max\bigl(2+3\alpha,\min(1+5\alpha,3+2\alpha)\bigr)$ для методов с узлами в пределах полутора шагов; численно установлены лишь существование методов с указанными показателями и устойчивая неудача поиска методов с большими. Начиная с $\alpha=\frac23$ проявляется дробный аналог глубоких барьеров Бутчера, и при $\alpha\to1$ дробный барьер непрерывно переходит в классическое ограничение $p\le4$. Без ограничения на узлы <<цепочечная>> ветвь формально продолжается до $\alpha^*\approx0.72$ за счёт методов с неограниченно растущими коэффициентами.
\item \emph{Численно} установлено, что вычисление лаг-члена по значениям стадий ограничивает глобальный порядок величиной $\approx1+\alpha$ (строгое доказательство этого наблюдения не проводилось); вычисление по значениям в узлах сетки с интерполяцией достаточной степени позволяет реализовать показатель $q$, что подтверждено прямым измерением локальной ошибки и расчётами на линейной, нелинейной и неавтономной системной задачах. Превышение прогноза на отдельных задачах объясняется малостью коэффициента при главном члене (для полной цепочки он пропорционален $\|J\|^m$) и не противоречит теории.
\end{enumerate}

Перспективным направлением дальнейших исследований является доказательство гипотез~\ref{conj4} и~\ref{conj5} и описание барьеров при произвольном числе стадий, в том числе строгий анализ случая разрыва цепочки $a_{i+1,i}=0$ и третьей ветви барьера. Отдельного внимания заслуживают стартовая процедура вблизи $t=0$, где решение негладкое (ряды по $t^{j+k\alpha}$~\cite{Lubich1983}), сгущающиеся сетки, неявные методы и вопросы устойчивости, а также строгая оценка влияния погрешности лаг-члена. Полученные результаты создают теоретическую основу для построения явных методов повышенного порядка точности для дробно-дифференциальных уравнений, возникающих при локализации задач наследственной механики~\cite{Prik2026de}, в том числе в составе сеточно-характеристических схем~\cite{PP2026zhvm}.

\section*{Благодарности}
Автор выражает глубокую благодарность научному руководителю, члену-корреспонденту РАН И.\,Б.~Петрову за внимание к работе и ценные обсуждения.

\section*{Финансирование работы}
Исследование выполнено за счёт гранта Российского научного фонда № 23-11-00035-П, \url{https://rscf.ru/project/23-11-00035-П/}.

\section*{Конфликт интересов}
Автор данной работы заявляет, что у него нет конфликта интересов.

\begin{thebibliography}{99}\small
\bibitem{Volterra1982} Вольтерра~В. Теория функционалов, интегральных и интегро-дифференциальных уравнений. М.: Глав. ред. физ.-мат. лит., 1982. 300~с.
\bibitem{Rabotnov1966} Работнов~Ю.\,Н. Ползучесть элементов конструкций. М.: Наука, 1966. 752~с.
\bibitem{Rabotnov1977} Работнов~Ю.\,Н. Элементы наследственной механики твёрдых тел. М.: Наука, 1977. 384~с.
\bibitem{SKM1987} Самко~С.\,Г., Килбас~А.\,А., Маричев~О.\,И. Интегралы и производные дробного порядка и некоторые их приложения. Минск: Наука и техника, 1987. 688~с.
\bibitem{Uchaikin2008} Учайкин~В.\,В. Метод дробных производных. Ульяновск: Артишок, 2008. 512~с.
\bibitem{Diethelm2010} Diethelm~K. The Analysis of Fractional Differential Equations. Berlin: Springer, 2010.
\bibitem{PPKh2024} Петров~И.\,Б., Приказчиков~Д.\,А., Хохлов~Н.\,И. Численное решение интегродифференциальных уравнений теории вязкоупругости с ядрами экспоненциального и работновского типов // Докл. РАН. Матем., информ., проц. упр. 2025. Т.~521. С.~88--95. DOI: 10.31857/S2686954325010116.
\bibitem{PP2026zhvm} Приказчиков~Д.\,А., Петров~И.\,Б. Численный метод решения уравнения теории линейной вязкоупругости в слоистых средах // Ж. вычисл. матем. и матем. физ. 2026 (в печати).
\bibitem{PP2025ljm} Prikazchikov~D.\,A., Petrov~I.\,B. Determination of parameters of integral operators for the initial boundary value problem of the theory of viscoelasticity on creep curves // Lobachevskii J. Math. 2025. V.~46. P.~3413--3419.
\bibitem{Prik2026sibzhim} Приказчиков~Д.\,А. Обобщение SOE-подхода на нелинейные операторы вязкоупругости с повреждениями // Сиб. журн. индустр. матем. 2026 (в печати).
\bibitem{Prik2026de} Приказчиков~Д.\,А. Локализация интегродифференциальных уравнений теории линейной вязкоупругости с ядром Работнова // Дифференц. уравнения (направлена в редакцию).
\bibitem{Ghaffari2023} Ghaffari~R., Ghoreishi~F., Saad~N. Fractional order Runge--Kutta methods // Fractal and Fractional. 2023. V.~7, №~3. Art.~245. arXiv:2210.13138.
\bibitem{Pouzet1960} Pouzet~P. Méthode d'intégration numérique des équations intégrales et intégro-différentielles du type de Volterra de seconde espèce. Formules de Runge--Kutta // Symposium on the Numerical Treatment of Ordinary Differential Equations, Integral and Integro-Differential Equations (Rome, 1960). Basel: Birkhäuser, 1960. P.~362--368.
\bibitem{Lubich1983} Lubich~Ch. Runge--Kutta theory for Volterra and Abel integral equations of the second kind // Math. Comp. 1983. V.~41, №~163. P.~87--102.
\bibitem{Lubich1985} Lubich~Ch. Fractional linear multistep methods for Abel--Volterra integral equations of the second kind // Math. Comp. 1985. V.~45, №~172. P.~463--469.
\bibitem{Garrappa2010} Garrappa~R. Order conditions for Volterra Runge--Kutta methods // Appl. Numer. Math. 2010. V.~60, №~5. P.~561--573.
\bibitem{BHN1982} Brunner~H., Hairer~E., N{\o}rsett~S.\,P. Runge--Kutta theory for Volterra integral equations of the second kind // Math. Comp. 1982. V.~39, №~159. P.~147--163.
\bibitem{LubichOstermann1993} Lubich~Ch., Ostermann~A. Runge--Kutta methods for parabolic equations and convolution quadrature // Math. Comp. 1993. V.~60, №~201. P.~105--131.
\bibitem{BanjaiLopez2018} Banjai~L., L\'opez-Fern\'andez~M. Efficient high order algorithms for fractional integrals and fractional differential equations // arXiv:1803.00642. 2018.
\bibitem{DFF2002} Diethelm~K., Ford~N.\,J., Freed~A.\,D. A predictor--corrector approach for the numerical solution of fractional differential equations // Nonlinear Dynamics. 2002. V.~29. P.~3--22.
\bibitem{Brunner2004} Brunner~H. Collocation Methods for Volterra Integral and Related Functional Equations. Cambridge: Cambridge University Press, 2004.
\bibitem{OdibatShawagfeh2007} Odibat~Z.\,M., Shawagfeh~N.\,T. Generalized Taylor's formula // Appl. Math. Comput. 2007. V.~186. P.~286--293.
\bibitem{Butcher2016} Butcher~J.\,C. Numerical Methods for Ordinary Differential Equations. 3rd ed. Chichester: Wiley, 2016.
\bibitem{HNW1993} Hairer~E., N{\o}rsett~S.\,P., Wanner~G. Solving Ordinary Differential Equations I. Nonstiff Problems. 2nd ed. Berlin: Springer, 1993.
\bibitem{PolyaSzego} P\'olya~G., Szeg\H{o}~G. Problems and Theorems in Analysis II. Berlin: Springer, 1976 (Part Five, Ch.~1).
\end{thebibliography}

\end{document}
